%\documentclass[apj]{emulateapj}
\documentclass[preprint2,twocolumn,times,tighten]{aastex63}

\usepackage{graphicx,times}
\usepackage{subfigure}
\newcommand{\be}{\begin{equation}}
\usepackage{threeparttable}
\usepackage{booktabs}
\newcommand{\ee}{\end{equation}}
\newcommand{\bea}{\begin{eqnarray}}
\newcommand{\eea}{\end{eqnarray}}
\usepackage{color}
\newcommand*{\al}{\color{black}}
\newcommand*{\xu}{\color{magenta}}


\usepackage{enumerate}
\usepackage{amsmath}
\usepackage{cases}
\usepackage{longtable}
\usepackage{hyperref}
\usepackage{epstopdf}
\usepackage{amsmath,bm}
\usepackage{amssymb}
\usepackage{natbib}
\usepackage{morefloats}
\usepackage{multirow}
\usepackage{array}
\usepackage{verbatim}
\usepackage{lineno}
%\linenumbers


\newcommand{\erf}{\mathrm{erf}\,}

\begin{document}




\title{Mirror acceleration of cosmic rays in a high-$\beta$ medium}
%{Propagation and Acceleration of Cosmic Rays: effects of magnetic field wandering and magnetic mirrors} 

\email{lazarian@astro.wisc.edu; sxu@ias.edu}


\author{A. Lazarian}
	\affiliation{Department of Astronomy, University of Wisconsin, 475 North Charter Street, Madison, WI 53706, USA}
	\affiliation{Centro de Investigación en Astronomía, Universidad Bernardo O’Higgins, Santiago, General Gana 1760, 8370993, Chile}
	\author[0000-0002-0458-7828]{Siyao Xu\footnote{NASA Hubble Fellow}}
	\affiliation{Institute for Advanced Study, 1 Einstein Drive, Princeton, NJ 08540, USA}
	

\begin{abstract}
In a weakly compressible high-$\beta$ medium,
pitch-angle scattering and the associated scattering acceleration of cosmic rays (CRs) by anisotropic 
Alfv\'{e}n and slow modes of 
magnetohydrodynamic (MHD) turbulence
is inefficient. 
To tap the energy from magnetic compressions for efficient particle acceleration, 
a diffusion mechanism that can effectively confine particles in space without causing their trapping 
or pitch-angle isotropization
is needed. 
We find that 
the mirror diffusion in MHD turbulence recently identified in 
\citet{LXm21} 
satisfies all the above conditions and serves as a promising diffusion mechanism for efficient acceleration of CRs via their stochastic 
non-resonant interactions with magnetic compressions/expansions.
The resulting mirror acceleration is dominated by the slow-mode eddies with their lifetime comparable to the mirror diffusion time of CRs.
Consequently, we find that 
the acceleration time of mirror acceleration
is independent of the spatial diffusion coefficient of CRs. 
The mirror acceleration brings new life for the particle acceleration in a weakly compressible/incompressible medium
and has important implications for studying CR re-acceleration in the high-$\beta$ intracluster medium. 


%We study how the acceleration of cosmic rays (CRs) propagating in realistic turbulent magnetic fields is affected by non-resonant large-scale magnetic compressions. We identify the cosmic ray superdiffusion perpendicular to the magnetic field direction as an important process that brings to life a new type of second-order Fermi acceleration, namely, the one arising from bouncing from magnetic mirrors. It is due to this type of perpendicular diffusion that CRs diffuse parallel to the magnetic field while bouncing from magnetic mirrors.  

\end{abstract}

\section{Introduction} 

Cosmic rays (CRs) are the non-thermal energetic component of galaxies and galaxy clusters.
Studies on their diffusion and acceleration have far-reaching astrophysical implications on understanding the strong correlation between nonthermal emission and infrared luminosities of star-forming galaxies 
\citep{Yun01,Ajel20}
and influence of CRs on star formation, and galaxy and galaxy cluster evolution
(e.g., \citealt{Blas00,Pfro08,Hop12,Evoli12,La16,BruL16,Wien17,Hol18,Krum20,Seme21,Quat22,Liu23}.)




The interaction of CRs with the ubiquitous turbulent magnetic fields causes their diffusion and stochastic acceleration
\citep{Schlickeiser02,Mar20,Liu21}.
The advances achieved recently in theoretically understanding
\citep{GS95,LV99}
, simulating
(e.g., \citealt{MG01,CL02_PRL,Bere14}), 
and observationally measuring  
(e.g., \citealt{Laz18,Humc19,Huse22,Yuen22})
magnetohydrodynamic (MHD) turbulence 
bring significant changes to the standard paradigm of CR diffusion and acceleration
(e.g., \citealt{Chan00,YL04,XY13,LY14,XLb18,LucS18,Siou20,Hucr22,Kemp22,Bea22,LXm21,Forn21,Zhda21,Natt21,Lemo22,Samp23,Lemo23,Kemp23}).


In a weakly compressible high-$\beta$ medium, where $\beta$ is the ratio of the gas pressure to the magnetic pressure, both Alfv\'{e}n and slow modes of MHD turbulence have scale-dependent anisotropy 
\citep{CV00,CL02_PRL}
and are thus inefficient in scattering the CRs with the Larmor radii much smaller than the energy injection scale of turbulence 
\citep{Chan00,YL02,BYL2011,XL20}. 
Consequently, the spatial confinement and 
stochastic acceleration via gyroresonant scattering is inefficient. 
In a high-$\beta$ medium, 
the resonance-broadened Transit Time Damping
(TTD) by slow modes is found to dominate over that by fast mode in scattering and accelerating CRs 
\citep{XLb18}. 
However, as the ``pitch-angle scattering" by TTD is in fact caused by the stochastic acceleration in the direction parallel to the magnetic field, TTD alone cannot spatially confine CRs. 
It would lead to a highly anisotropic pitch angle distribution and cannot be self-sustained. 




\begin{comment}
Charged energetic particles or cosmic rays present an important ingredient of physics of space science and 
astrophysical processes. It is a custom to use the term "energetic particles" in space physics, which includes the 
processes acceleration and propagation in the Solar atmosphere, solar wind, Earth magnetosphere and heliosphere. The energetic particles outside the our direct neighborhood, i.e. of galactic and extragalactic origin are usually termed cosmic rays. As a rule, the latter are discussed in terms of particles of higher energies. However, in terms of the physics of processes that we discuss in this paper, there is no difference between the aforementioned types of particles 
and, for the sake of simplicity, we will not distinguish of the energetic particles of different origins and will use the term "cosmic rays" for them.

"Cosmic rays" produced by solar activity are important for the predictions of the effect of the temporal variations flux of these particles on sensitive space-based equipment, transmission of radio waves in the Earth upper atmosphere, etc., and is the subject that is of the primary interest of the "space weather" research direction (see....). The cosmic rays in our galaxy have pressure comparable to pressure arising from interstellar gas and magnetic field and therefore they are dynamically important. 

Cosmic rays interact strongly with magnetic fields and their motion perpendicular to magnetic field lines is limited. The interaction of cosmic rays with magnetic field is a subject of intensive research which analytically is usually done within the framework of the so-called Quasilinear Thieory (QLT) (see Jokipii1966, Melrose..., Schlikeiser 2006). The discrepancies of the theoretical expectations and numerical expectations has become clear with the progress of numerical simulations (see Qin et al. 2002, see the review by Gabici et al. 2019).

Understanding of cosmic ray propagation and the corresponding interaction with turbulent magnetic fields are important for the quantitative description of key astrophysical and Space science processes, e.g. for predicting the arrival of the energetic particles from Solar flares to the Earth, explaining the chemical composition of cosmic rays, accounting for the ionization of gas in molecular clouds and circumstellar accretion disks, launching galactic winds as well as returning the particles into the zones of acceleration within the shock and reconnection acceleration  processes. Thus the subject of cosmic ray propagation has been extensively explored. 

The progress in describing the dynamics of cosmic rays became possible with the advent of numerically tested theory of MHD turbulence (Goldreich \& Sridhar 1995; Lazarian \& Vishniac 1999; Cho \& Vishniac 2000;
Maron \& Goldreich 2001; Cho et al. 2002; Cho \& Lazarian 2003; Kowal et al. 2012; Beresnyak 2014, see also book by Beresnyak \& Lazarian 2019), which has also received support from the observations of magnetic fluctuations the solar wind (Horbury et al. 2008; Forman et al. 2011). The calculations of cosmic ray scattering and acceleration made with this theory (Chandran 2000b; Yan \& Lazarian 2004; Brunetti \& Lazarian 2007; Yan \& Lazarian 2008; Xu \& Yan 2013; Lazarian \& Yan 2014; Xu et al. 2016; Xu \& Lazarian 2018; Sioulas et al. 2020) got very different results from the calculations based on the adhoc models of MHD turbulence, e.g. on the basis of the empirical model with adjustable parameters, the so-called slab and 2D model (see Matthaeus et al. 1990, Bieber .... , monograph by Schalchi...). These efforts allowed to provide solutions to some long-standing problems of cosmic ray research, e.g. to reconcile the amplitude of the magnetic fluctuations observed in the solar wind with the relative inefficiency of cosmic ray scattering by these fluctuations, to resolve the puzzle related to the so-called 90 degree scattering (see Yan \& Lazarian 2008), provide the quantitative description of the propagation of cosmic rays perpendicular to the mean magnetic field (see Lazarian \& Yan 2014). 
\end{comment}


Non-resonant interactions of CRs with %weakly compressible 
MHD turbulence have also been investigated for studying CR diffusion and stochastic acceleration
\citep{Noe68,CesK73,Ptuskin88,Klep95,CL06,Brunetti_Laz,Med15,BruL16,XZ17,Bres22,LXm21}.
The acceleration of CRs by magnetic compressions on scales larger than their Larmor radii was studied by 
\citet{CL06},
%can cause the increase of particle momentum in both parallel and perpendicular directions with respect to the magnetic field. 
%The acceleration efficiency is sensitive to the diffusion of CRs along the magnetic field, i.e., parallel diffusion. 
where
as the diffusion due to pitch-angle scattering was invoked, 
significant compression of gas is required for 
a net energy gain in average
(see also \citealt{Druf12}).
Therefore, their model cannot lead to efficient acceleration in a weakly compressible medium. 



Based on the modern understanding of MHD turbulence, 
\citet{LXm21} recently identified the mirror diffusion of CRs parallel to the magnetic field. 
Unlike the mirror trapping 
with particles trapped between two magnetic mirror points, 
the mirror diffusion happens due to the perpendicular superdiffusion of turbulent magnetic fields 
\citep{LVC04,XY13,Bere13,Eyin13,LY14,Hucr22}.
CRs interact with magnetic compressions, i.e., magnetic mirrors, along the magnetic field lines
while experiencing superdiffusion in the direction perpendicular to the magnetic field. 
In a high-$\beta$ medium, as the pitch-angle scattering is inefficient, the mirror diffusion is expected to be the dominant mechanism for parallel diffusion
\citep{LXm21}.
\footnote{We note that other mechanisms in a high-$\beta$ medium can also contribute to confining CRs, e.g., field line wandering in super-Alfv\'{e}n turbulence 
\citep{Brunetti_Laz}, strong scattering by small-scale sharp magnetic field bends
\citep{Lemo23,Kemp23}.}
Unlike pitch-angle scattering, 
mirroring can spatially confine particles without causing stochastic change in their pitch angles. 
Therefore, it provides a promising diffusion mechanism for the non-resonant acceleration by large-scale\footnote{By ``large scale", we mean the scales larger than the particle Larmor radius.} magnetic compressions to be efficient in a weakly compressible medium. 





%At the same time, scattering is not the only process that affects the propagation of cosmic rays in magnetized media. Cosmic ray trapping by large-scale magnetic irregularities was earlier studied by a number of authors, e.g., Fermi (1949); Noerdlinger (1968); Cesarsky \& Kulsrud (1973); Klepach \& Ptuskin (1995). However, these efforts were not connected to the framework of a realistic model of turbulence.  

%The importance of the effect of trapping on propagation was demonstrated in the realistic model of MHD turbulence in Xu \& Lazarian (2020, henceforth XL20). In terms of physics, the model in XL20 is similar to that in Cesarsky \& Kulsrud (1973), i.e. the role of trapping is limited to constraining the range of pitch angles over which the diffusion induced by scattering takes place. The destiny of the trapped particles was addressed in Lazarian \& Xu (2021, henceforth LX21 ). There it was demonstrated that, due to the perpendicular superdiffusion (Lazarian \& Yan 2014), the particles between the magnetic mirrors are not trapped, but involved in the diffusion, which was termed "bouncing diffusion". The nature of this new effect that we also discuss in \S.. is related to low statistical probability of CRs to retract their trajectories as they follow rapidly diverging field lines (LV99, Yan \& Lazarian 2008, Eyink et al. 2011, Xu \& Yan 2013).  As the result the CRs exhibit superdiffusive behavior and bounce from different magnetic mirrors. As a result, LX21 presented the actual CR diffusion parallel to magnetic field as a sum of two types of diffusion, the classical scattering one described in XL20 and the new bouncing diffusion along the mean magnetic field.  


In this work, by taking into account the mirror diffusion, we will reexamine the non-resonant acceleration of CRs by magnetic compressions in a high-$\beta$ medium, which we will term ``mirror acceleration".
In Section 2, we will review the mirror diffusion in a high-$\beta$ medium.
In Section 3, we will focus on the mirror acceleration in a high-$\beta$ medium. 
Conclusions are presented in Section 4. 

%This paper explores the effects of the aforementioned bouncing diffusion on the second order Fermi Acceleration. This new type of acceleration, which we term "bouncing acceleration" would not be present if the CRs were trapped. It is the "untrapping" of CR due to the perpendicular superdiffusion that brings to life this new type of acceleration.





\begin{comment}
\section{Basic properties of MHD turbulence}

\subsection{Basic scalings for Alfvenic turbulence}

MHD turbulence is a major agent determining the dynamics of CRs (Schlikeiser 2006). The turbulence in magnetized media in typical astrophysical settings is being injected at the scale $L_i$ at the injection velocity $V_L$ and the energy then cascades to smaller scales. If the injection happens with $V_L>V_A$, where $V_A$ is the Alfven velocity, this is the case of super-Alfvenic turbulence. If $V_L<V_A$ this is the case of subAlfvenic turbulence. The ratio $M_A=V_L/V_A$ is the Alfven Mach number. $M_A=1$ corresponds to the transAlfvenic turbulence.

 As CRs have Larmor radii significantly larger than the thermal ions and therefore the small scale physics of the MHD cascade is irrelevant. Note, that turbulence assumes the cascading and therefore Alfven waves that are created by CR instabilities, e.g. by streaming instability (see Farmer \& Goldreich 2003), gyro-resonance instability (see Lazarian \& Beresnyak 2006) are not classified as turbulence. 

If turbulence is injected at $M_A>1$ the magnetic field is of marginal importance at the injection scale $L_i$ and it cascades initially in a hydrodynamic turbulence way, i.e. is essentially, Kolmogorov turbulence. However, as the turbulent velocities decrease as $v_l\sim V_L (l/L_i)^{1/3}$ the role of magnetic force increases. Eventually, at the scale
\begin{equation}
l_A=L_i M_A^{-3}
\label{lA}
\end{equation}
the turbulent velocity gets equal to $V_A$ and the turbulent becomes fully magnetohydrodynamic. 

The theory of transAlfvenic MHD turbulence was developed by Goldreich \& Sridhar (1995, henceforth GS95). The subsequently developed theory of turbulent reconnection (Lazarian \& Vishniac 1999, henceforth LV99) provides a way both to explain GS95 theory in analogy with the Kolmogorov theory and, very importantly, augment the GS95 theory with an essential concept of "local system of reference". Indeed, LV99 predicted that magnetic reconnection happens in one eddy turnover time. Therefore the fluid motions that mix up magnetic fields without bending them, i.e. eddies with axes aligned with the ambient magnetic field, should present the prevalent motions. The cascade of such eddies is going to be Kolmogorov-type with
\begin{equation}
v_l\sim V_A \left(\frac{l_\bot}{L_i}\right)^{1/3},
\label{GS95perp}
\end{equation}
where it is taken into account that for the transAlfvenic turbulence the injection velocity is equal to $V_A$ and $l_\bot$ is the scale of the eddies perpendicular to the magnetic field. In this picture it is obvious that the magnetic field at hand is the {\it local} magnetic field that permits the eddy. This direction of the local magnetic field can differ substantially from the mean magnetic field direction, but it is the only field that the eddy feels and it is the field that determines the unconstrained eddy rotation. Note, that the original GS95 theory did not use the local magnetic field concept and all the calculations were performed in terms of the mean magnetic field. However, numerical simulations in Cho \& Vishniac (2000), Maron \& Goldreich (2001), Cho et al. (2002) confirmed that the universal scalings between the parallel and perpendicular sizes of the eddies, i.e. $l_\bot$ and $l_\|$ exists only in the frame of local magnetic field and is not valid if the parallel and perpendicular motions are measured in respect to the mean magnetic field. 

Within the aforementioned model of eddies aligned with the local magnetic field it is easy to obtain the relation between $l_\bot$ and $l_\|$. Indeed, consider an eddy with the perpendicular size of $l_\bot$. Turbulent reconnection allows its unconstrained turnover about the direction of the ambient magnetic field. This motion, however, induces an Alfvenic perturbation that propagates with the speed $V_A$. The timescale of this perturbation $l_\|/V_A$ should be equal to the eddy turnover time $l_\bot/v_l$. The corresponding relation
\begin{equation}
\frac{l_\|}{V_A}\approx \frac{l_\bot}{v_l},
\label{critbal}
\end{equation}
was termed "critical balance" in GS95. However, we would like to note that in GS95 the original meaning of $l_\|$ and $l_\bot$ is different, as those quantities are measured in terms of mean magnetic field direction. Using the Kolmogorov scaling of the eddies given by Eq. (\ref{GS95perp}) and the critical balance relation given by Eq. (\ref{critbal}) one can obtain the scale-dependent anisotropy of Alfven modes:
\begin{equation}
l_{\|}\approx L_i \left(\frac{l_\bot}{L}\right)^{2/3}
\label{ll}
\end{equation}
for the $M_A=1$ setting. Note that Eqs.  (\ref{GS95perp}) and (\ref{ll}) should be understood in the statistical sense. They represent the relations between the most probable values quantities involved. For instance, on the basis of numerical simulations Cho et al. (2002) provided an analytical fit for distribution function of $l_\bot$ for a given $l_\|$. This was used later in Yan \& Lazarian (2002, 2004) and subsequent studies for CR scattering.

Returning back to $M_A>1$ MHD turbulence, it is evident that $l_A$ given by Eq. (\ref{lA}) should be used in Eqs. (\ref{GS95perp}) and (\ref{ll}) instead of $L_i$ in order to describe the MHD cascade at scales less than $l_A$ (Lazarian 2006). The corresponding equations written in terms of $L_i$ include the factors $M_A$ and $M_A^{-1}$, respectively.

The situation is more involved in case of subAlfvenic turbulence, i.e. for $M_A<1$. It was shown in LV99 that at the injection scale the turbulence is {\it weak} and it is only the perpendicular scale of the motions that is changing, i.e. $l_\|=L_i$. The scaling obtained in LV99 for the weak turbulence in the assumption of the isotropic turbulent driving at scale $L_i$ is:
\begin{equation}
v_l\approx V_L \left(\frac{l}{L_i}\right)^{1/2}
\end{equation}
and this result was supported in the subsequent studies by Galtier et al. (2001). 

The intensity of interactions of Alfvenic cascade is increasing as $l_\bot$ decreases. It was shown in LV99 that at a scale
\begin{equation}
l_{trans}\approx L_i M_A^2,
\label{ltrans}
\end{equation} 
where $M_A<1$, the turbulence gets strong. For the sub-Alfvenic MHD turbulence LV99 provided the relations 
\begin{equation}
v_l\approx V_L (l/L_{i})^{2/3} M_A^{4/3},
\label{vAlf}
\end{equation}
and 
\begin{equation}
l_{\|}\approx L_i \left(\frac{l_\bot}{L_i}\right)^{2/3} M_A^{-4/3},
\label{lpar}
\end{equation}
which for $M_A=1$ transfer into our previous expressions for transAlfvenic turbulence given by Eqs. (\ref{GS95perp}) and (\ref{ll}).


\subsection{Compressible fluctuations in MHD turbulence}

We are interested in the magnetic bottles created by compressible MHD turbulence. In this case,
the scaling magnetic perturbations change with the size of the perturbation is determined by compressible
motions in the MHD cascade. These motions can be described by the slow and fast modes. We use word
"modes" rather than waves, as the cascading is making the properties of the fluctuations different from
the properties of linear waves. In particular, it was shown numerically in Cho \& Lazarian (2003) that the slow
modes in both low- and high-$\beta$
plasmas are sheared by Alfven modes and therefore get the 
Goldreich \& Sridhar (1995, henceforth GS95) scaling. This scaling corresponds to the Kolmogorov-type $k_{\bot}^{-5/3}$ spectrum in terms of perturbations perpendicular to magnetic field and $k_{\|}^{-2}$ spectrum of perturbation in the direction parallel to magnetic field. Here it is necessary to explain that the $\bot$ and $\|$ should be understood in terms of {\it local} direction of magnetic field, rather than the mean field. This distinction was not initially a part of MHD turbulence as it suggested in GS95. It is, however, evident in the picture of magnetic eddies that arises within the picture of fast turbulent reconnection in Lazarian \& Vishniac (1999, henceforth LV99). Indeed, the eddies rotating perpendicular to the local magnetic field that percolates the eddies do not feel the resistance from the field, provided that the magnetic field rearrangement, i.e. magnetic reconnection, happens within one eddy turnover time. This is exactly the prediction of LV99 theory. The alignment of magnetic field in respect to the local magnetic field direction is confirmed in numerical studies (Cho \& Vishniac 2000, Maron \& Goldreich 2001, Cho et al. 2002). The theoretical discussion of scaling of slow modes in Lithwick \& Goldreich (2001), as well as further theoretical studies, (e.g. Xu et al. 2016) are done in terms fluctuations defined in respect to the local direction of magnetic field.  

From the discussion above it follows that slow modes of MHD turbulence follow $k_{\bot}^{-5/3}$ and $k_{\|}^{-2}$
spectrum, which due to the aforementioned requirement of the local of magnetic field cannot be measured in terms
of the global mean magnetic field. Fortunately for the problem at hand, namely, for the propagation and acceleration of
energetic particles by MHD turbulence, the particles do feel exactly the local magnetic field that they gyrate around. This changes a lot in terms of the description of the energetic particle dynamics. 

For instance,  the standard derivations
of the quasi-linear theory (QLT) is formally only true for infinitesimally weak turbulent perturbations, the case that is usually of little practical interest as the fluctuations at the injection scale can be comparable or larger than the mean field. 

Nevertheless, in the local system of reference, the particles interact with the small scale perturbations, while the large scale, e.g. magnetic fluctuations at the injection scale, can be arbitrary. Naturally, this extends significantly the QLT applicability.

The scaling of fast modes is somewhat less certain. The theoretical considerations in GS95 in terms of high-$\beta$ plasma and in Cho \& Lazarian (2002, 2003) in terms of low-$\beta$ plasma suggested the "isotropic" spectrum of fast modes with $k^{-3/2}$. The "isotropic" should be understood in a sense that $E(k)\sim ...$, which is different from the accepted expression for the truly isotropic spectrum of a vector field in hydrodynamics (see Monin \& Yaglom 1972). The differences are discussed in detail in Lazarian \& Pogosyan (2012). 

The uncertainty of the actual spectrum of the fast modes arises from the fact that later higher resolution numerical calculations in Kowal \& Lazarian (2010) indicated that the measured spectrum is close to $k^{-2}$ rather than $-3/2$. This discrepancy may arise from the differences of the codes used in Cho \& Lazarian (2003) and Kowal \& Lazarian (2010) or it can be also similar to the transition from acoustic-type turbulence with $k^{-3/2}$ to the turbulence of shocks with $k^{-2}$ spectral slope.  

When treating perturbations arising from MHD turbulence, in particular, from slow and fast modes that are essential for
cosmic ray trapping, it is natural to identify the magnetic perturbations $\delta b$ with the scale-dependent magnetic perturbations of the corresponding modes, i.e. with $\delta b_k$. The use the actual scaling of the slow and fast modes makes our study directly applicable to predicting cosmic ray trapping in the realistic astrophysical setting.  




\subsection{Scaling of slow and fast MHD modes}

Magnetic compressions produced by fast and slow modes produce magnetic mirrors which properties we explore in this study. 

\subsubsection{Low beta case}

The relation between the amplitudes of Alfven, slow and fast MHD modes may depend on the properties of driving.
Assuming that the driving is happening with 
$\delta V_A=\delta V_s=\delta V_f$, Cho \& Lazarian (2003)
obtained the following relations for magnetic compressions induced by the modes for low $\beta$:
\begin{equation}
  \frac{\delta B}{B_0} \|_s\sim \beta M_s,   
\end{equation}
and 
\begin{equation}
  \frac{\delta B}{B_0}\|_f\sim \beta^{1/2} M_s.   
\end{equation}

We note, that this is a linear approximation and in case of supersonic motions it will fail. In particular, we expect to observe shocks for $M_s>1$. Those can induce compression $\sim M_s^2$ increasing the contribution of slow modes in magnetic field compression.


It was discussed in GS95 that when Alfvenic turbulence is in the strong regime the slow modes are slaved by it.  In that regime, Alfven modes evolve fast within one period and shear slow modes, inducing their cascading. However, the GS95, is formulated for transAlfvenic turbulence, i.e. for $M_A=1$ regime.

Consider the cascade of slow modes for $M_A<1$. In this situation, the 
 the Alfvenic turbulence is in the weak regime for a range of scales, from the injection scale $L$ to the transition scale $l_{tr}=LM_A^2$. Over these scales, the shearing of the slow modes by Alfvenic perturbations  proceeds in the perpendicular direction due to the weak cascade. Using the classical expressions of the non-linear Alfven wave interactions (Iroshnikov 1963, Kraichnan 1965) one gets the non-linear cascading rate of Alfvenic turbulence
\begin{equation}
\tau_{nl}^{-1} \approx \frac{(kv_k)^2}{\omega_A}
\label{Krai}
\end{equation}
where the angular frequency $\omega_A=k_\| V_A$. 

Eq. (\ref{Krai}) can be rewritten in terms of involved scales as
\begin{equation}
\tau_{nl}^{-1}\approx \frac{l_\| v_l^2}{l_{\bot}^2 V_A}
\end{equation}.

As it was explained in LV99 (see Appendix A there) that the problem of the earlier studies was that the Alfvenic turbulence anisotropy was not accounted for.  However, in reality, the cascade is very anisotropic. In the weak regime, it involves only the changes in the perpendicular scale, while keeping $l_\|=const$. Assuming that the injection of energy is performed isotropically at the scale $L$, and taking into accout the constancy of the energy cascade $\epsilon = v_l/\tau_{nl}$ one gets for the scaling of the Alfvenic modes in weak turbulence  
\begin{equation}
    v_l\approx V_A M_A \left(\frac{l_\bot}{L}\right)^{1/2}
\end{equation}
 As $l_\bot=l_{tr}$ the turbulence transfers to the strong regime with the critical balance between given by Eq. (\ref{critbal}). The cascading rate for the strong turbulence is equal to 
 \begin{equation}
     \tau_{nl}^{-1}\approx \frac{V_A}{l_\|}
     \label{casc_strong}
 \end{equation}

 

At the same time, the slow modes are cascading in parallel direction through interacting with other slow modes. In low-$\beta$ case the slow modes propagate with sound velocity $c_s$ and their interactions in accoustic-like cascade changing $l_{\|}$. Eq. (\ref{Krai}) changes in this case to accommodate this type of interactions.  Substituting $\omega_{sl}=c_s/l_\|$ instead of $\omega_A$ in Eq. (\ref{Krai}), results in the non-linear interaction rate 
\begin{equation}
\tau_{nl, s}^{-1} \approx \frac{(v_{l,s}/l_\|)^2}{c_s/l_\|}=\frac{v_{l,s}^2}{l_\| c_s}
\label{sl_nonl}
\end{equation}
Substituting this into the turbulence cascading rate
\begin{equation}
\epsilon \approx \frac{v_{l,s}^2}{\tau_{nl,s}}=\frac{V_{inj, s}^4}{L c_s}
\label{epsslow}
\end{equation}
we get the velocity scaling for the change of the slow mode with the parallel scale
\begin{equation}
v_{l,s}\approx c_s M_s \left(\frac{l_\|}{L}\right)^{1/4}~~~M_s<1.
\label{vsl}
\end{equation}
Combining this with Eq. (\ref{epsslow}) one gets 
\begin{equation}
\tau_{nl, s}^{-1}\approx \frac{V_A}{L} M_A M_s \left(\frac{L}{l_\|}\right)^{1/2}~~~~M_s<1
\label{nonlinearS}
\end{equation}
where the Alfven and sonic Mach numbers $M_A$ and $M_s$ are defined in terms of the slow mode injection velocity. 




For sufficiently large $M_s$, slow modes can create shocks. Then the scaling would change to 
\begin{equation}
    v_{l,s}\approx c_s M_s \left(\frac{l_\|}{L}\right)^{1/2},~~~M_s>1
    \label{casc_Ms>1}
\end{equation}
which provides the cascading rate
\begin{equation}
    \tau^{-1}_{nl,s}\approx \frac{V_{inj} M_s}{L}.
    \label{casc_M>1}~~~M_s>1.
\end{equation}

If the rates given by Eq. (\ref{casc_Ms>1}) and Eq. (\ref{nonlinearS}) compared with rates of strong turbulence cascading given by Eq. (\ref{casc_strong}, it is evident that only in the case of $M_s>1$ the self-cascading by slow modes prevail. The slow modes have their self-generated structure over the range of scales $[L/(M_A M_s), L]$, while the on the scales less than $L/(M_A M_s)$ the slow modes are still slaved by Alfvenic cascade and mimic the structure of Alfvenic mode. 

For the range of slow mode dominance not to be empty, the product $M_A M_s$ should be larger than one. This can be achieved if plasma $\beta$ parameter is less than $M_A^4$, which limits the applicability of the effect.

The parallel structures arising from the accoustic-like cascade of slow modes can be important for special circumstances. One of them is when the damping of Alfvenic turbulence prevents its transition to strong regime, i.e. $d_{damp,\bot}>l_{tr}$. In this case, the perpendicular structures are dominated by Alfvenic turbulence, while parallel ones are due to slow mode cascade. The structure of slow modes in this situation is slab turbulence reminiscent to the 2D plus slab models sometimes employed in space science research (see Mathaeus...).
This can happen with both $M_s>1$ and $M_s<1$, but the requirement $d_{damp,\bot}>l_{tr}$ is very restrictive. In addition, the fast modes should be disregarded to adopt the 2D plus slab model. 

The structure of turbulence reminiscent of 2D plus slab can also emerge in the presence of a Alfvenic turbulence over a limited range. To get this effect, $L/(M_A M_s)>d_{damp,\|}$, where $d_{damp,\|}$ is the parallel damping scale for the Alfvenic cascade. In this case, on the over the range from $L$ to
\begin{equation}
    l_\bot \approx l_{tr} (d_{damp,\|}/L_i)^{3/2},
    \label{edge}
\end{equation}
the Alfvenic cascade is strong, but does not impose its scaling of the slab structure of the slow modes. For scales less than given by Eq. (\ref{edge}) the perpendicular shearing by Alfven modes, i.e. "2D" modes are absent, while the "slab" created by slow modes proceeds to the slow mode dissipation scale. 


In the case of $M_A M_s>1$, the cascading slow modes create additional parallel structures that can increase scattering and the second order Fermi acceleration of CRs. 

The fast modes in low beta medium are magnetic field compression. They create their own cascade according to Cho \& Lazarian (2002). Their interaction with Alfven modes is reduced due to the difference in cascading rate and with slow modes by the disparity of the velocities. The accoustic cascade of fast is isotropic.    

\subsection{High $\beta$ case}

The compression arising from slow and fast modes in high $\beta$ medium for the incompressible driving in the linear approximation  are (Cho \& Lazarian 2003):
\begin{equation}
    \frac{\delta B_s}{B_0} \sim M_A,
\end{equation}
and
\begin{equation}
    \frac{\delta B_f}{B_0}\sim  M_s^2,
\end{equation}
with the ratio of the two compression equal to $\beta/M_A$. Therefore, the magnetic compression by fast modes is subdominant for nearly incompressible media, i.e. for $\beta\gg 1$.

Without Alfvenic modes, slow modes can produce the cascade with
\begin{equation}
    v_l\approx V_A M_A \left( \frac{l_\|}{L_i}\right)^{1/4}.
\end{equation}
This induces the non-linear cascading rate
\begin{equation}
    \tau^{-1}_{nl}\approx \frac{V_A}{L_i} M_A^2\left(\frac{L_i}{l_\|}\right)^{1/2}, 
\end{equation}
which for $M_A<1$ that we consider cannot be larger than the parallel cascading rate of the Alfven cascade given by Eq. (\ref{casc_strong}). 

Therefore, the slow modes can preserve their cascade only in the case that the parallel damping scale of Alfvenic cascade is larger than $L_i$.








 





 

 

\section{CRs motion perpendicular to magnetic field }

\subsection{Magnetic field meandering}

Magnetic field makes the propagation of CRs anisotropic. For instance, the CR diffusion parallel and perpendicular to magnetic field is different. If we disregard magnetic field wondering and the perpendicular transposition by the CR Larmor radius $r_L$ takes place as a CR moves a mean free path $\lambda\gg r_L$ along magnetic field. Thus, the motion in the perpendicular direction corresponds to the random walk with the mean free path of $r_L$. This difference in the parallel and perpendicular diffusion resulted in the paradox. The expected propagation of CR perpendicular to the galactic magnetic field should be strongly reduced, which was at odds with the observational data. 

The solution for the paradox was proposed by Jokipii (...) who noticed that the magnetic field meandering dominates the perpendicular CR propagation. In the absence of the theory describing MHD turbulence and the meandering, the quantitative description of this process was impossible. This resulted in a number of well-entrenched fallacies related to CR dynamics.
For instance,  the existence of diffusive propagation along magnetic field lines raised a possibility that the perpendicular transport can be subdiffusive provided that the CRs retrace their paths along magnetic field lines. Indeed, if magnetic field meandering obeys the traditional random walk pattern and there is a diffusive transport along magnetic field lines the transposition in the direction perpendicular to magnetic field direction will proceed subdiffusively with $\langle R^2_\bot \rangle \sim t^{1/2}$ (Kota \& Jokipii...).  

The situation changed when quantitative description of magnetic field meandering in 3D strong MHD turbulence was presented in Lazarian \& Vishniac (1999, henceforth LV99). There the magnetic field meandering was described as a part of the process of 3D magnetic reconnection. However, the application of the obtained laws to the relevant fields was straightforward and was done in later papers. In particular, it was applied to heat transfer in plasmas (Narayan \& Medvedev 2002, Lazarian 2006) and cosmic ray dynamics (Yan \& Lazarian 2008, Lazarian \& Yan 2014), correcting a number of misconceptions existed in the literature. 

With the LV99 predictions of magnetic meandering confirmed numerically (Lazarian et al. 2004, Beresnyak 2013, Xu \& Yan 13), the theory provides a solid basis for the description of the perpendicular propagation of CRs. In what follows we describe the theory and provide its modifications that the scattering of particles entails.

The dynamics of the magnetic field are different if the magnetic field is strong, i.e., magnetic energy is equal to or larger than the kinetic energy associated with the driving. In this situation, the ability of magnetic field line stretching is constrained, and the kinematic dynamo does not operate.

For the magnetic field separation, it is important to distinguish the motions parallel and perpendicular to the locally defined mean magnetic field. The strong magnetic field prevents magnetic field stretching but, in the presence of fast reconnection,\footnote{The latter can be shown to be true for the turbulent magnetic field (LV99), but for the sake of discussion, we do not dwell on this issue at this point.} does not prevent the mixing motions of the magnetic field perpendicular to the mean magnetic field. 

In the presence of turbulence, the one-scale driving is applicable to the motions less than the turbulence damping scale, assuming that the Ohmic resistivity is less than viscosity, i.e., that the fluid has the Prandtl number larger than unity. In this situation, the turbulent cascade stops at the scale of critically damped eddies $l_{min,\bot}$, and the magnetic field lines at the scales less than this scale are subject to the shear from such eddies. The eddies in strongly magnetized turbulence are anisotropic. Thus the critically damped eddies have dimensions $l_{min,\bot}$ and $l_{min,\|}$ with $l_{min,\|}>l_{min,\bot}$. In this setting, Eq. (\ref{3}) should be modified to account for the anisotropy of shear, i.e.
\begin{equation}
    dy=\frac{y}{l_{min,\|}} ds,
    \label{6}
\end{equation}
the latter equation reflects the fact that the shear along the magnetic field is determined by the size of the eddies along the local magnetic field $l_{min,\|}$.

One can understand the meaning of different terms in Eq. (\ref{6}) by noticing that $y/l_{min,\bot}$ is the tangent of the angle that the magnetic field line has with the direction of the mean magnetic field. Thus multiplying this tangent by $ds$, one gets the increment of $dy$ perpendicular to mean magnetic field as the magnetic field line is traced over $ds$. 

Naturally, the solution of Eq. (\ref{6}) is modified compared to Eq. (\ref{5}), namely, 
\begin{equation}
s_{RR,\bot}=l_{min,\|}\ln \frac{l_{min,\bot}}{y_{0, \bot}},
\label{7}
\end{equation}
where $y_{0,\bot}$ should be at least larger than the thermal electron Larmor radius. If Eq. (\ref{7}) is used to trace CRs, then the minimal value of $y_{0,\bot}$ corresponds to the CR Larmor radius. 

For turbulence within the inertial range, the perpendicular dispersion of magnetic field lines is determined by the shear at the scale of the magnetic separation. The tangent of the inclination of magnetic field lines for an eddy with $l_{\bot}$ and $l_\|$ is
$l_\bot/l\|$. Therefore the perpendicular to mean magnetic field transposition of a magnetic field line is given by
\begin{equation}
    dy=\frac{l_\bot}{l_\|} ds.
    \label{8}
\end{equation}

For strong MHD turbulence, one can use the critical balance condition and rewrite as
\begin{equation}
    dy=\frac{v_l}{V_A}ds,
    \label{9}
\end{equation}
while for the regime of weak turbulence, the approximation of constant shear is not applicable due to the oscillating nature of turbulence at scales larger than $l_{tr}=LM_A^2$. 

The solution of Eq. (\ref{9}) can be obtained by substituting the scaling of strong turbulence for $v_l$. The corresponding solution is 
\begin{equation}
    y^2\approx \frac{s^3}{L}M_A^4,
\end{equation}
which coincides with the expression for the magnetic field line dispersion first obtained in LV99, but in a different way. 

The present derivation unifies the concepts of Lyapunov growth of the magnetic field line separation and magnetic field wandering. It shows that magnetic stochasticity can be present and important even without turbulence. However, turbulence significantly accelerates the stochasticity by increasing the shear that magnetic fields experience. 




\subsection{Meandering and ballistic propagation of CRs}

Alfvenic turbulence dominate the meandering/wandering of magnetic field lines (see LV99). This induces the CRs dispersion perpendicular to magnetic field, as CRs propagate ballistically along magnetic field lines. Naturally, this dispersion will change with Alfven Mach number $M_A$. 

According to LV99, the MHD turbulence is a collection of eddies with their axes aligned along the direction of the nearby magnetic field. Therefore, it is natural to assume that the magnetic field lines deviate over the perpendicular scale of the eddy $l_{\bot}$ while tracing the magnetic field over the parallel scale of the eddy $l_\|$. The additions of $l_\bot$ can be both positive and negative and therefore it is the dispersion $\langle y^2 \rangle$ that is increasing in the random walk fashion with the distance measure along the magnetic field line $s$:
\begin{equation} 
d\langle y^2 \rangle\approx l_\bot^2 \frac{ds}{l_\|}, 
\label{y2}
\end{equation}
where $l_\|$ and $l_\bot$ are related through
Eq. (\ref{lpar}).

The peculiar feature of the random walk of $y^2$ is that the step of the random walk increases with $s$. Indeed, consider dispersion magnetic field lines that at $s=0$ are at the distance $d_{diss, \bot}$ from each other. The initial separation of the lines is governed by the dynamics of the eddies with $l_\bot\approx d_{diss,\bot}$, but as magnetic field lines separate further, the dominant contribution comes from larger eddies, i.e. with $l_\bot=y>d_{diss,\bot}$. This entails the accelerated separation of the field lines. 
In view of the considerations above, 
it is natural to identify $l_\bot$ in Eq. (
\ref{y2}) with $\langle y^2\rangle^{1/2}$. Within this approximation one can obtain
\begin{equation}
   \langle y^2\rangle^{1/3} \approx \frac{1}{3}\left(\frac{M_A^4}{L}\right)^{1/3} + C
\end{equation},
where the constant $C$ in our setting is equal to initial separation $d^{2/3}$ that should be equal or larger than $d_{diss,\bot}^{2/3}$. In many practical settings the initial separation can be disregarded and then one can get the LV99 expression 
\begin{equation}
\langle y^2\rangle \approx \alpha \frac{s^3}{L_i} M_A^4, ~~~s<L_i,
\label{perpen}
\end{equation}
where $\alpha$ is a prefactor which exact value is difficult to establish from the simple scaling arguments above. Instead, it should be obtained empirically from numerical simulations. 

The considerations above assume that the initial separation of magnetic field lines is larger than the dissipation scale of turbulence. The minimal separation of magnetic field is given by the Larmor radius of the electric current carrier, i.e. $r_e$. In this case, $r_e< d_{diss\bot}$ the magnetic field lines initially do not experience the effect of turbulent cascade. Instead, they are being moved apart by the shear that arises from the marginally damped eddies at the scale $d_{diss,\bot}$. As a result, the magnetic field lines initially deviate exponentially and they deviate by the distance $d_{diss,\bot}$ over the parallel scale $s_{RR}$ that is calculated in Appendix \ref{sec:small_r} and is shown to be $\approx d_{diss,\bot}^{2/3} L_i^{1/3} M_A^{-4/3} \ln(d_{diss,\bot}/r_e)$, where the prefactor is the scale of the parallel dissipation scale $d_{diss,\|}$. In other words, before the closest magnetic fields experience superbalistic growth of dispersion, they have to move some distance $s_{RR}$ in a different regime. However, due to the properties of the logarithm function, $s_{RR}$ is not very different from $d_{diss,\|}$. The difference decreases further if one considers not the perpendicular dispersion of the closest possible magnetic field lines, but the perpendicular dispertion of cosmic rays. Then instead the minimal separation of CRs is not $r_e$ but the Larmor radius of a CRs $r_0\gg r_e$. In fact, for sufficiently energetic CRs, $r_0$ is larger than $d_{diss,\bot}$ and such CRs cannot experience anything but superbalistic dispersion. All in all, in most cases, the presence of the finite dissipation scale $d_{diss,\bot}$ is important only for relatively low energy CRs and over a limited range $s_{RR}$ which is a few times of $\sim d_{diss,\|}$. 


The dispersion of magnetic field lines is different over the scales much larger than $L_i$. In this situation, the dispersion is growing in the direction perpendicular to magnetic field in a random walk fashion with the random walk step equal to the $l_{trans}$ given by Eq. (\ref{ltrans}). Therefore, the corresponding perpendicular dispersion of magnetic field lines is 
\begin{equation}
 \langle y^2\rangle \sim (L_i s)^{1/2} M_A^2, 
\label{eq:random}
\end{equation}
(LV99). The cosmic rays with the mean free path $\lambda\gg L_i$ reproduce this dispersion. They get perpendicular displacement of $l_{tr}$ for every distance $L_i$ that the cover along magnetic field lines.
Further we discuss the consequences of the magnetic field dispersion for CRs propagation when the CRs mean free path $\lambda$ is less than $L_i$.



\subsection{Perpendicular diffusion in the presence of trapping and scattering}
\label{s:perp&scat}


The perpendicular meandering of magnetic field lines in the random walk fashion given by Eq. (\ref{eq:random}) does not induce subdiffusion, provided that the $d_{diss,\bot}<l_{tr}$. CRs moving the distance $L_i$ experience by scattering, but they do not retrace their path along the same magnetic field lines. The superdiffusion of magnetic field lines on scales less than $l_{tr}$ induces CRs follow new magnetic field lines (Yan \& Lazarian 2008, Lazarian \& Yan 2014). The subdiffusion is possible only in a  setting with $d_{diss,\bot}>l_{tr}$. Another case of subdiffusion can take place over parallel scales less than $s_{RR}$.

Trapping of CR is the effect recently discussed in Xu \& Lazarian (2020). The CRs are reflected from the magnetic mirrors, which modifies their parallel propagation. This also changes perpendicular to magnetic field propagation of CRs. 

The law given by Eq. (\ref{perpen}) is applicable to the perpendicular transposition of CR $y_{cr}$ when free ballistic motion of particles is considered (see Lazarian \& Yan 2014). In this situation $s=v_\|t$, where $v_\|$ is the parallel to magnetic field velocity of CR and $t$ is the time of CR travel. This true as the particle motion reflects the superdiffusion of magnetic field lines. 

For the case of CRs diffusing along magnetic field lines, 
$s^2\sim t$. The resulting 
 perpendicular motion is still superdiffusive in time.
i.e. 
$\sqrt{\langle y_{cr}^2 \rangle} \sim t^{3/4}$ (see Yan \& Lazarian 2008, Xu \& Yan 2013, Hu et al. 2021), but the growth rate of CR perpendicular dispersion is reduced due to reduced rate of parallel propagation. The scaling is somewhat different in terms of the relations between the parallel and perpendicular transposition.

 In the presence of strong scattering, the arguments above can be modified to describe the perpendicular superdiffusion of CRs while they propagate diffusively along the field lines. 

For the parallel propagation one has
\begin{equation}
    s^2\approx D_\| t\approx\lambda_{\|} v_{cr} t,
\end{equation}
where $v_{cr}$ is the velocity of CRs parallel to the magnetic field direction, and $D_\|$ is the parallel diffusion coefficient. 

The  angle of a given magnetic field line with respect to the mean magnetic field can be approximated as
\begin{equation}
    \alpha (s)\approx \frac{\langle y^2 (s) \rangle^{1/2}}{s}\approx \frac{s^{1/2}}{L^{1/2}} M_A^2,
    \label{alpha}
\end{equation}
where the scaling of perpendicular magnetic field line superdiffusion is used.
It is easy to see that in the  system of reference of the mean magnetic field, the perpendicular displacement of CR $\lambda_{cr\bot}\approx \alpha (s) \lambda\|$ if the CR moves over the mean free path $\lambda_\|$ along the magnetic field. The CR has a component of velocity  perpendicular to the mean magnetic field  $v_{cr, \bot}\approx \alpha (s) v_{cr, \|}$. 

The perpendicular propagation is a random walk with the step size and the velocity that increase with increasing with $s$. The parallel diffusion coefficient is $D_\|\approx v_{cr, \|} \lambda_\|$. Therefore for a given $s$ the perpendicular diffusion is $\sim v_{cr \bot} \lambda_\bot\sim \alpha^2 (s) D_\|$. The diffusion for larger $s$ dominates the resulting diffusion. Thus, the propagation of perpendicular to the mean magnetic field can be approximated as 
\begin{equation}
\langle y^2_{cr} \rangle \approx \alpha^2 D_\| t\approx \frac{s^{3}}{L} M_A^4,~~~ ~~M_A<1,~~\lambda_\|\ll s, 
\label{perp_diff}
\end{equation}



 
 However, the particles that experience scattering do the random walk steps both in parallel and perpendicular directions. If a particle moves a distance $l_\|$ before reversing its direction it also moves $l_\| \cos\phi$ in the perpendicular direction, where $\phi$ is the angle of local magnetic field in respect to the mean field. After the reversing the direction the particle can move in the same direction or also reverse the direction of its perpendicular displacement. Therefore, this is a typical case of the random walk, which on the scales much larger than the injection scale is expected to result in the relation between the parallel $D_{\|}$ and perpendicular $D_{diff, \perp}$ diffusion coefficients of particle (Yan \& Lazarian 2008, Lazarian \& Yan 2014):
 \begin{equation}
D_{diff, \perp}\approx M_A^4 D_\|.
\label{eq: diff}
\end{equation}

The perpendicular diffusion, consider CR covering the distance $\beth\gg LM_A^2$ in the perpendicular direction. According to Eq. (\ref{diff2}) the magnetic field lines that guide the diffusion undergo a random walk. As the magnetic field lines is traced over $L$ in parallel direction, its perpendicular transposition is $LM_A^2$. In the limit of CR Larmor radius $r_L$ much smaller than the cosmic ray mean free path $\lambda$, diffusing CRs follow closely wandering magnetic field lines. To cover the distance $\beth$ with the elementary step $LM_A^2$, $N=(\beth/LM_a^2)^2$ steps are required. The time for each step is $\delta t_{step}\approx L^2/D_\|$. The total time to cover the distance $\beth$ is $\delta t \approx N \delta t_{step}$. Inserting this into the expression for the diffusion coefficient one gets
\begin{equation}
    D_{\bot}\approx \frac{\beth^2}{\delta t}\approx D_\| M_A^4,
    \label{diffperpnew}
\end{equation}
which coincides with the CR diffusion coefficient in Yan \& Lazarian (2008). The $M_A^4$ dependence was confirmed in  Xu \& Yan 2013 through numerical simulations.


 
 \subsection{Perpendicular diffusion for subAlfvenic turbulence}

If turbulence is injected isotropically at scale $L$ with $V_{inj}$ that is less than the Alfven velocity $V_A$, there is a range of scales $[L_i, l_A]$, where $l_A=LM_A^2$ for which the turbulence is weak (LV99, Lazarian 2006). The scaling of the weak Alfvenic turbulence is $E(k)\sim k^{-2}$ (LV99, Gaultier et al. 2000) is different from the spectrum of the strong GS95 turbulence $E(k)\sim k^{-5/3}$. More importantly, the weak cascade increases only the perpendicular wavenumber, while leaving the parallel wavenumber $k_{\|}$ constant and equal to $\sim 1/L_i$.  If the CRs diffuse along magnetic field with the diffusion coefficient $D_\|$, they cover the distance $L$ over the time $t_{step}\approx L_i^2/D_\|$. During this step time the perpendicular transposition of CRs that follow magnetic field lines
is $l_A$. During time $t$ the perpendicular displacement is  $\sim l_A^2 t/t_{step}$, which provides the total perpendicular diffusion coefficient given by Eq.(\ref{eq: diff}).  
 
\end{comment}  


\section{Mirror diffusion in a high-$\beta$ medium}

\subsection{Interaction of particles with magnetic mirrors}


For a particle with the total momentum $p$ moving along the magnetic field of strength $B$, the first adiabatic invariant is conserved during the mirroring process,
\begin{equation}
    J_1 \propto p^2\frac{1-\mu^2}{B},
    \label{J1}
\end{equation}
where the cosine of the pitch angle $\mu$ changes in accordance with the change of $B$, 
and pitch angle is the angle between the directions of ${\bm p}$ and ${\bm B}$. 

As the first condition 
for the mirroring to happen, 
the Larmor radius $r_g$ of the particle should be smaller than the size of the mirror, i.e., the variation scale of the magnetic field. 
%A particle traveling in a magnetic inhomogeneity with a scale much larger than the 
%particle Larmor radius preserves its adiabatic invariant, i.e. $p_{\bot}^2/B=const$, provided that the scale of the inhomogeneity is much larger than the particle Larmor radius $r_L$. 
As a particle gyrates moving along a magnetic field of increasing strength, its perpendicular
momentum $p_\perp = p \sqrt{1-\mu^2}$
increases (see Eq. \eqref{J1}). 
With a constant particle energy, it follows that 
the parallel momentum $p_\| = p \mu$
of the particle decreases. 
%If the magnetic field increases to the value $B_0+\delta b$,
%where $B_0$ is the initial magnetic field, the particle perpendicular momentum $p_{\bot}$ gets equal to the total moment
%of the particle $p$. As a result, the particle's parallel momentum vanishes, and the particle bounces back, experiencing the gradient of the magnetic field. The magnetic field inhomogeneity acts as a spring that first slow the particle and then reflects it back. The difference from the spring, however, is that the energy for bouncing is not being accumulated in the spring, but is stored in the particle component of perpendicular momentum $p_\bot$.
When the magnetic fluctuation is sufficiently large, the mirroring force is so strong that $p_\|$ can decrease to zero. 
Then the particle is reflected by the mirror. 
The condition for the mirror reflection is  
\begin{equation}
\frac{p_\bot^2}{B_0}=\frac{p^2}{B_0+\delta b},
\end{equation}
where $B_0$ and $B_0 + \delta b$ are the magnetic field strengths in the weak and strong magnetic field regions.
The maximum $\mu$ corresponding to the smallest pitch angle for mirror reflection is  
%Introducing the cosine of the pitch angle, the ratio of $p_{\bot}/p$ one gets the condition for the pitch angle critical for bouncing:
\begin{equation}
\mu^2_{max}=\frac{\delta b}{B_0+\delta b}.
\label{mu}
\end{equation}
%All particles, irrespective of their energy, will bounce back from the magnetic perturbation $\delta b$ if their pitch angle $\mu$ is less than $\mu_c$.  At the same time, particles with $\mu>\mu_c$ escape the magnetic bottle. Thus $\mu_c$ also defines the escape cone of particles. 
The particles with $\mu < \mu_{max}$ are subject to mirroring. 

%The value of $\mu_c$ is being modified in the presence of resonance scattering (Kulsrud ...). In particular, it was shown in XL20 that fast modes are most important decreasing the value of $\mu_c$. This effect is essential for understanding of magnetic field diffusion, where $\mu_c$ serves as a borderline between the particles undergoing classical diffusion determined by scattering and those which undergo a new type of diffusion introduced in LX21. The latter process termed "bouncing diffusion" is typically much slower than the scattering diffusion, which has many important astrophysical implications (Xu 2021...). The scattering provides redistribution between the two populations of particles, which is a very important effect for particle propagation (LX21, Xu 2021) On the contrary, in this paper we focus on the acceleration of the particles undergoing bouncind diffusion. Therefore, for the sake of simplicity, we do not consider resonant scattering and our critical pitch angle is given by Eq. (\ref{mu}). 

As the third condition for mirroring, 
the pitch-angle scattering rate cannot exceed the mirroring rate for $J_1$ to be conserved. 
In a low-$\beta$ medium, with efficient scattering by fast modes, mirroring happens over a range of $\mu$ with the mirroring rate larger than the scattering rate
\citep{CesK73,XL20,LXm21}.
For the weakly compressible high-$\beta$ medium considered in this work, 
fast modes are sound waves moving with an infinite velocity and their effect on particle dynamics is negligible. 
The scattering by Alfv\'{e}n and slow modes is inefficient at all $\mu$'s, 
and thus the mirroring occurs within the entire range 
$0<\mu<\mu_{max}$. 











\begin{comment}
\subsection{Role of scattering}

In the presence of scattering, the critical value of $\mu$ changes, as scattering allows trapping particles to escape magnetic bottles. The scattering by fast modes was identified in Yan \& Lazarian (2002) as the principal mechanism for particle scattering in the environments where the injection scale $L$ is much larger than the particle Larmor radius $r_L$. Thus in what follows, we focus on fast mode scattering.



The rate of change in $\mu$ due to scattering as the scattering rate according to XL20 is
\begin{equation}\label{eq: scraft}
    \Gamma_{s,f} = \frac{2 D_{\mu\mu,\text{QLT},f}}{\mu^2} \approx
          \frac{ \pi}{28} \frac{ \delta B_f^2}{B_0^2}   \Big(\frac{v}{L\Omega}\Big)^\frac{1}{2}
                \Omega  (1 - \mu^2) \mu^{-\frac{3}{2}}  .
\end{equation}
where the approximate expression of fast mode pitch angle diffusion coefficient $D_{\mu\mu,\text{QLT},f}$ 
\citep{Xuc16,XLb18}
\begin{equation}\label{eq: appfqltu}
\begin{aligned}
      D_{\mu\mu,\text{QLT},f} 
      %&= \frac{\pi^2}{2} C_\mu C_f \frac{1}{v_\|} \int  d\eta dk \eta (1-\frac{x^2}{4})  k^{-\frac{3}{2}}  \delta(   k  -  \frac{\Omega}{v_\| \eta} ) \\
       %                        & = \frac{\pi^2}{2} C_\mu C_f  \Omega^{-\frac{3}{2}} v_\|^{\frac{1}{2}} \int  d\eta   (1-\frac{ (1-\mu^2) (1-\eta^2) }{4 \mu^2 \eta^2})  
        %                        \eta^{\frac{5}{2}} \\
       %                        & = \pi^2 C_\mu C_f  \Omega^{-\frac{3}{2}} v_\|^{\frac{1}{2}}
         %                       \Bigg[\int  d\eta \eta^{\frac{5}{2}} -  \frac{1-\mu^2}{4 \mu^2}\int d\eta   (1-\eta^2)  
          %                      \eta^{\frac{1}{2}}  \Bigg] \\
           %                    & = \pi^2 C_\mu C_f  \Omega^{-\frac{3}{2}} v_\|^{\frac{1}{2}}
           %                       \Bigg[ \frac{2}{7} -  \frac{2}{21} \frac{1-\mu^2}{ \mu^2}     \Bigg] 
       %  &  \approx \frac{2 \pi^2}{7}  C_f  \frac{\Omega^\frac{1}{2}}{B_0^2}   v^{\frac{1}{2}}  (1 - \mu^2) \mu^{\frac{1}{2}} \\
         &  \approx \frac{ \pi}{56} \frac{ \delta B_f^2}{B_0^2}   \Big(\frac{v}{L\Omega}\Big)^\frac{1}{2}
                \Omega  (1 - \mu^2) \mu^{\frac{1}{2}}  .
%         &  = A  \frac{ \delta B_f^2}{B_0^2}   (1 - \mu^2) \mu^{\frac{1}{2}} 
\end{aligned}
\end{equation}
was used.

This rate should be compared to the rate of $\mu$ change induced by effect of the mirror.

The rate of the adiabatic change in $\mu$ due to the spatial variation of magnetic field is
(\citealt{CesK73},XL20)
\begin{equation}\label{eq: gtraft}
\begin{aligned}
   \Gamma_{b,f}  =   \Big |\frac{1}{\mu} \frac{d\mu}{dt} \Big| &=  \frac{v}{2B_0} \frac{1-\mu^2}{\mu}  b_{fk} k  \\
 & = \frac{v}{2L} \frac{\delta B_f^4}{B_0^{4} } \frac{1-\mu^2}{\mu^7}   
\end{aligned}
\end{equation}
 used  the magnetic perturbation of fast modes $b_{fk}$ a
\begin{equation}\label{eq: fascc}
    b_{fk} = \delta B_f (kL)^{-\frac{1}{4}}
\end{equation}
according to the scaling of isotropic fast modes 
\citep{CL02_PRL}.

Eq.(\ref{eq: gtraft}) is valid for sufficiently small Larmor radii of particles , i.e. $\mu>\mu_\text{min,f}$ 
\begin{equation}
     \mu_\text{min,f} = \sqrt{\frac{b_{fk} (r_L)}{B_0}}  = \sqrt{\frac{\delta B_f}{B_0}} \Big( \frac{r_L}{L} \Big)^{\frac{1}{8}},
\end{equation}
and 
$b_{fk} (r_L)$ is the magnetic perturbation at $r_L$.

For $\mu < \mu_\text{min,f}$, the bouncing rate is 
\begin{equation}\label{eq: trfasmu}
\begin{aligned}
     \Gamma_{b,f}  & =   \frac{v}{2B_0} \frac{1-\mu^2}{\mu}  b_{fk}(r_g) r_g^{-1} \\
                                                                  & =  \frac{v}{2 r_g } \frac{ \delta B_f}{B_0} \Big( \frac{r_g}{L} \Big)^{\frac{1}{4}}  \frac{1-\mu^2}{\mu} .
\end{aligned}
\end{equation}


By equating the rate of $\mu$ change between scattering and bouncing, i.e., $\Gamma_{s,f} = \Gamma_{b,f}$, XL20 obtained 
(Eqs. \eqref{eq: scraft} and \eqref{eq: gtraft}, 
\begin{equation}\label{eq: faucf}
    \mu_{c,scat1} \approx  \bigg[ \frac{14}{\pi} \frac{ \delta B_f^2}{B_0^2} \Big(\frac{v}{ L \Omega}\Big)^\frac{1}{2} \bigg]^\frac{2}{11}, 
\end{equation}
which defines the critical cosine of pitch angle in the presence of fast mode scattering. 

In the regime when the finite Larmor radius of particles is important, one can obtain:
\begin{equation}
    \mu_{c,scat2} \approx
\end{equation}
which presents an important case for bouncing of particles with large $r_L$. 

In what follows, we will also use 
\begin{equation}
    \mu_{c, scat}=\max [\mu_{c, scat1}, \mu_{c, scat2}],
    \label{eq:mucscat}
\end{equation}
to provide the distinction of bouncing and non-bouncing particles.

\end{comment}

\subsection{Mirror diffusion of energetic particles}

\subsubsection{Mirror trapping and mirror diffusion}

{By following \citet{CesK73}}, 
in \citet{XL20}, the interaction of particles with MHD turbulence was 
separated into two regimes, {the classical diffusion caused by pitch-angle scattering} at $\mu>\mu_{c}$ and mirror trapping of particles at $\mu<\mu_{c}$. 
At the critical $\mu_c$, there is a balance between mirroring rate and scattering rate
\citep{XL20,LXm21}. 
{Within a limited range of pitch angles for scattering to dominate over mirroring, 
the corresponding scattering mean free path is reduced, and the problem of vanishing scattering close to $90^\circ$, i.e., $90^\circ$ problem
\citep{Jokipii1966}, 
is naturally resolved by mirroring. }
%The trapping decreased particles' parallel diffusion and helped address the problem of 90-degree scattering.\footnote{The pitch angle diffusion coefficient $D_{\mu \mu}$ arising from resonance scattering vanishes as $\mu \rightarrow 0$. In the classical treatment, this should result in catastrophic accumulation of particles with $\mu=0$. Trapping of with $\mu<\mu_{c, scat}$ means that to reverse the particle direction $D_{\mu\mu}$ should not be zero at $\mu_{c,scat}\neq 0$ which is true for pitch angle scattering. The reversal of particle direction is accomplished for $\mu<\mu_{c,scat}$ through bouncing.}

{In \citet{LXm21} (hereafter, LX21), the dynamics of mirroring particles at 
$\mu<\mu_{c}$ 
was reexamined. 
There it was found that due to the perpendicular superdiffusion of magnetic field lines in Alfv\'{e}nic turbulence
\citep{LV99,Eyin13,Laz20}, 
the mirroring particles can not be trapped between two mirroring points 
as they can not trace the same magnetic field line back and forth. 
Instead, they are stochastically reflected by different magnetic mirrors
while moving along the diverging turbulent magnetic field lines.
Thus a new diffusion process termed ``mirror diffusion" was identified. 
%The bouncing diffusion can arise in MHD turbulence from slow and fast modes. The rate of the bouncing diffusion was shown in LX21 to be much slower than the scattering diffusion of non-bouncing particles with $\mu>\mu_{c, scat}$. 
It is found to be much slower than the scattering diffusion
\citep{LXm21,Xu21}. }
% focused on the stochastic acceleration of bouncing particles.

Figure \ref{fig: sket} illustrates the process of mirror diffusion. The mirroring particles do not retrace their way back but follow new routes {led by stochastic magnetic field lines.} 
In a high-$\beta$ medium, the perpendicular superdiffusion induced by Alfv\'{e}n modes 
and the mirroring induced by slow modes together result in the mirror diffusion of CRs. 



\begin{figure}[htbp]
%\centering   
%\subfigure[]{
%   \includegraphics[width=7cm]{skea.png}\label{fig: sketa}}
%\subfigure[]{
   \includegraphics[width=7cm]{skeb.png}
\caption{ Illustration for mirror diffusion.
The perpendicular superdiffusion of {\it turbulent} magnetic field lines originally identified in \citet{LV99}, induces the parallel diffusion of particles that {stochastically interact with different} magnetic mirrors.  
Black lines are magnetic field lines. Blue and red lines are trajectories of two particles with a small initial separation.  From LX21. }
\label{fig: sket}
\end{figure}

%As we showed earlier, the perpendicular to the magnetic field superdiffusion of particles obeys Eq. (\ref{eq:bouncesuper}). This superdiffusion makes the retracing of particle trajectories improbable. Sampling different eddies and bouncing from the magnetic mirrors associated with the eddies, the particles diffuse in the parallel to magnetic field direction, which constitutes the parallel diffusion, i.e. bouncing diffusion, that is identified and quantified in LX21.


\subsection{Scattering diffusion and mirror diffusion}
\label{ssec:scatvsmirror}

Scattering diffusion is accompanied by the stochastic change of $\mu$. 
When a particle interacts with an eddy, 
if the scattering is so efficient that   
the timescale for pitch-angle isotropization is shorter than both the particle diffusion timescale 
and the lifetime of the eddy, 
the particle 
can simultaneously sample the compression (expansion) in the direction perpendicular to the magnetic field 
and the expansion (compression) in the direction parallel to the magnetic field, which accompany each other in a weakly compressible medium. 
After averaging over an isotropic distribution, 
energy gain and loss cancel out, and 
there is no net energy change. 
If the scattering is inefficient, 
particles freely move through the turbulent medium without being affected by magnetic compression/expansion, 
leading to inefficient acceleration. 

By contrast, 
mirror diffusion does not cause stochastic pitch angle change. 
A particle at a given $\mu$ can preferentially sample the compression/expansion in either perpendicular or parallel direction with respect to the magnetic field. 
Without averaging over an isotropic distribution, a net energy change
after interacting with an eddy is expected. 
Mirroring can also significantly suppresses the parallel diffusion of particles, which entails their sufficient interactions with magnetic compression/expansion. 
Therefore, mirror diffusion serves as a promising diffusion mechanism for efficient acceleration by large-scale magnetic compressions in a weakly compressible medium. 








\subsubsection{Mirror diffusion induced by slow modes}
\label{f}

In a high-$\beta$ medium, 
slow modes of MHD turbulence create magnetic compressions that {serve as magnetic mirrors}. 
For simplicity, we assume that the slow mode fluctuations corresponding to the parallel wavenumber $k_{\|,sl}$, i.e. $\delta b=b_{sk}$, are small compared to the mean magnetic field strength $B_0$. Thus Eq. (\ref{mu}) is approximately 
\begin{equation}
\mu_{max}^2\approx \frac{b_{sk}}{B_0}.
\end{equation}
%which gives the minimal magnetic fluctuation $b_{sk}\approx B_0 \mu^2$ for the mirroring of a particle with $\mu$. 
A particle with $\mu$ is most effectively reflected by the mirror 
with the magnetic fluctuation $b_{sk}$
corresponding to \citep{CesK73}
\begin{equation}
\mu^2\approx \frac{b_{sk}}{B_0}.
\label{mubk}
\end{equation}
%Note that $b_k$ will also trap particles with smaller $\mu$.
The scaling of slow modes is
\citep{CL03} 
%for a given $\mu$ one can define the minimal scale fluctuations that can trap the CRs.
%Indeed, for slow modes (see Cho \& Lazarian 2003)
\begin{equation}
\label{bslow}
   b_{sk} \approx  \delta B_s (k_{\|, sl} L)^{-\frac{1}{2}},
\end{equation}
where $\delta B_s$ is the amplitude of slow mode fluctuations at the injection scale $L$. Therefore, combining Eqs. (\ref{mubk}) and (\ref{bslow}) leads to
\begin{equation}
k_{\|, sl} \approx \frac{1}{L} \left(\frac{\delta B_s}{B_0 \mu^2}\right)^2 =\frac{1}{L\mu^4} \aleph_s^{2},
\label{bkslow}
\end{equation}
where 
\begin{equation}
    \aleph_s =  \frac{\delta B_s}{B_0}.
\end{equation}
For weakly compressible MHD turbulence, the equivalent expression for $\aleph_s$ is
\begin{equation}\label{equns}
    \aleph_s \approx \frac{V_{L,slow}}{V_A},
\end{equation}
where $V_{L,slow}$ is {the turbulent velocity of slow modes} at $L$. With this notation, the critical cosine of pitch angle for mirroring is 
\begin{equation}
    \mu_{c}\approx \aleph_s^{1/2},
    \label{mualeph}
\end{equation}
%{\xu provided $\aleph_s \ll 1$.}
which is the maximum $\mu$ determined by the largest magnetic fluctuation of slow modes. 
Particles with $\mu>\mu_c$ cannot be reflected through mirroring. 

The length corresponding to 
Eq. (\ref{bkslow})  
\begin{equation}
l_\mu \sim 1/k_{\|,sl},
\label{lmu}
\end{equation}
which acts as a proxy of the parallel mean free path of mirror diffusion at a given $\mu$ (LX21). 
%The fact that $l_\mu\sim \mu^4$ resolves the problem well-known problem of crisis of scattering at 90 degree angle. While resonance scattering in its classical treatment (see Jokipii...) is very inefficient in such scattering, the bouncing is the most efficient for such a setting.
{Its minimum value is determined by} 
\begin{equation}
 l_{\mu,min} = \text{max} [r_g, l_{\|,d}],
\end{equation}
where $l_{\|,d}$ is the parallel dissipation scale of slow modes. 
%In the presence of turbulence dissipation scale $l_{\|,min}$, the requirement is that $l_\mu\le l_{\|, min}$. 
The particles with $\mu$ less than  
(Eq. \eqref{bkslow})
\begin{equation}
    \mu_{min}\approx \aleph_s^{1/2}\left(\frac{l_{\mu, min}}{L}\right)^{1/4}
\end{equation}
%interact with scales larger with $l_{\| min}$. 
have the parallel mean free path equal to $l_{\mu,min}$.
Therefore, the parallel diffusion coefficient corresponding to mirror diffusion is 
\begin{equation}
    D_\| (\mu) \approx  v_p \mu l_\mu\approx v_p L \aleph_s^{-2} \mu^5  
    \label{Dpar}
\end{equation}
for $\mu > \mu_{min}$, 
where $v_p$ is the particle speed, and 
\begin{equation}
    D_\| (\mu) \approx  v_p \mu l_{\mu, min}
    \label{Dpar2}
\end{equation}
for $\mu < \mu_{min}$.



%\subsubsection{Mirror diffusion induced by fast modes}

 %Fast modes also induce bouncing diffusion. For $\mu < \mu_{c, scat}$, the bouncing effect regulates the parallel diffusion of particles with $\mu$. 
%{\xu In compressible MHD turbulence, fast modes can play a dominant role in mirroring due to its isotropic energy distribution 
%(LX21).}
%For particles with 
%\begin{equation}
%   \mu^2 \approx \frac{b_{fk}}{B_0}, 
%\end{equation}
%by using the scaling of fast modes in Eq. \eqref{eq: fascc},
%the parallel mean free path of mirror diffusion is 
%bouncing particles diffuse parallel to the magnetic field with a step size 
%\begin{equation}\label{eq: mirsz}
%    l_\mu = k_f^{-1} \approx L \Big(\frac{\delta B_f}{B_0}\Big)^{-4}  \mu^{8}
     %~~\mu_\text{min,f} <\mu <\mu_c ,
%\end{equation}
%for $l_\mu >l_{\mu,min}= \text{max}[r_g,l_d]$, or equivalently $\mu > \mu_\text{min}$.
%For $\mu < \mu_\text{min}$, the parallel mean free path is $l_{\mu,min}$.
%\begin{equation}
%   r_g = \frac{v_\perp}{\Omega} = \frac{v \sqrt{1-\mu^2}}{\Omega},
%\end{equation}
%The maximal $l_\mu$ corresponds to $\mu_{c}$,
%\begin{equation}\label{eq:lambdac}
%     l_{\mu,max}=L \Big(\frac{\delta B_f}{B_0}\Big)^{-4}  \mu_{c}^{8}.
%\end{equation} 
%The parallel diffusion coefficient is (LX21)
%\begin{subnumcases}
%     { D_{\|} (\mu)=\label{eq: dufup}}
%        v_p\mu l_\mu \approx v_p L \Big(\frac{\delta B_f}{B_0}\Big)^{-4}  \mu^{9},  \nonumber\\
%        ~~~~~~~~~~~~~~~~~~~~~~~~~~~~~\mu_\text{min} <\mu <\mu_{c} ,\\
%        v_p\mu l_{\mu,min}, ~~~~~~~~~~~~~~~~~~~~ \mu < \mu_\text{min},
%\end{subnumcases}
%which demonstrates a strong dependence on $\mu$.



%Similar to the case of slow modes, the turbulence dissipation at scale $l_{\|, min,f}$ modifies the diffusion at small $\mu$. Starting with 
%\begin{equation}
%    \mu_{diss, f} \approx
%\end{equation}
% particles can interact only with scales larger than $l_{\|, min,f}$.





\section{Mirror acceleration by slow modes in a high-$\beta$ medium}

%In a high-$\beta$ medium, 
%where $\beta$ is the ratio of gaseous to magnetic pressures, 
%the magnetic fluctuations that cause mirroring of particles arise from slow modes. 
%To focus on the effect of slow modes below, we consider MHD turbulence in {\xu a nearly incompressible/high-$\beta$} medium, where fast modes are sound waves moving with an infinite velocity and their effect on particle dynamics is negligible. 
In this section, we disregard scattering of particles by slow and Alfv\'{e}n modes for simplicity, 
which are inefficient in 
scattering particles with 
$r_g \ll L$
due to the turbulence anisotropy
\citep{Chan00,YL02,XL20}. 
%Therefore, for our analysis in this section, $\mu_c$ given by Eq. (\ref{mualeph}). 
We consider only the second-order Fermi/stochastic acceleration of mirroing particles with $\mu<\mu_c$. The stochastic acceleration of scattering particles was analysed in \citet{CL06}. 




\subsection{Electric field of slow modes}

{In addition to gyroresonant scattering, 
stochastic acceleration is frequently associated with 
%obligatory media compressions 
non-resonant interaction of particles with large-scale magnetic compressions
(see \citealt{CL06}).} 
 The latter can occur in both compressible and incompressible media. 
 %but is accompanied by fast diffusion in terms of pitch angle. 
In what follows, we will discuss the {non-resonant mirror acceleration} by slow modes in a weakly compressible medium. 
% This process results in an anisotropic distribution of accelerated particles. 
 
 %The magnetic field can be compressed in incompressible media by slow modes. However, this case requires a more careful approach. 
 In a weakly compressible media, the compression in one direction entails the expansion in {another} direction to 
 preserve the constant volume. 
 %preserve the volume of the magnetic flux tube constant. 
 %This invokes the separate treatment of the components of the CR momentum. 
{Accordingly, we will separately treat the perpendicular and parallel components of CR momentum.} 
 

 %To describe the process of CR acceleration in incompressible 
 Here we follow the approach  in 
\citet{CL06}. 
 Figure \ref{acceler} 
 illustrates the compression/expansion of the magnetic field and the acceleration/deceleration of a particle in a slow mode eddy. 
 The change of the magnetic field strength induces the electric field that is perpendicular to the magnetic field following the Faraday's law, 
 %which for the incompressible medium provides:
 \begin{equation}
     E=-\frac{1}{2\pi r_g c}\int ({\bf B}_0 \cdot \bigtriangledown v_l) d{\bf s},
     \label{E1}
 \end{equation}
where $d{\bf s}$ is the infinitesimal area element vector, 
and $c$ is the light speed.
 \begin{figure}
\centering
\includegraphics[width=0.9\linewidth]{sket.png}
\caption{The slow mode eddy is elongated along the local magnetic field direction, with a CR particle spiraling about the magnetic field and the electric field induced by the compression/expansion of the magnetic field. 
%The eddy extends along the magnetic field, improving the interaction of a CR and the eddy. 
Modified from \citet{CL06}.}
\label{acceler}
\end{figure}
For longitudinal slow modes, 
the velocity perturbation along the magnetic field induces magnetic compression/expansion perpendicular to the magnetic field. 
So the spatial derivative of velocity in Eq. (\ref{E1}) can be approximated by $v_{l,\|}/l_{\|}$, 
where {$v_{l,\|}$ is the parallel turbulent velocity at the parallel length scale $l_\|$}. 
It represents the rate of compression or expansion of the magnetic field. 
As a result, taking the integral over the particle orbit area $\pi r_g^2$, it is easy to obtain
\begin{equation}
    E\approx -\frac{B_0 r_g v_{l,\|}}{2l_\| c} 
\end{equation}
By substituting the expression of $r_g=p_\bot c/(q B_0)$, where $q$ is the electric charge, we can rewrite the above equation as  
 \begin{equation}
E\approx -\frac{v_{l, \|} p_\bot}{2 q l_{\|}}.
\end{equation}

The electric field that a {particle} experiences results in the change of its $p_\perp$:
\begin{equation}
 \frac{dp_\bot}{dt} =q E \approx -\frac{p_\bot v_{l,\|}}{2l_{\|}},
\label{pper}
\end{equation}
{where we 
%consider that $q$ is positive,}
%which results in changes in the kinetic energy associated with the perpendicular particle motion.
adopt 
%a convention that perpendicular to magnetic field compression corresponds to the 
{a negative sign for $v_{l,\|}/l_\|$ corresponding to magnetic compression 
and a positive sign corresponding to magnetic expansion.}. 
%i.e., $\frac{d p_\bot}{dt}>0$.
Naturally, strengthening of the magnetic field causes increase of $p_\perp$ based on the conservation of the first adiabatic invariant. 
%corresponds to $v_{l,\|}/l >0$. 
%{\bf AL  We should explain that the absolute value of $p_\bot$ cannot depend on the sign of the particle's charge. The perpendicular compression corresponds to the increase of $p_\bot$ for both electrons and protons. The momentum will be directed differently in space, but this does not matter. The increase of the absolute value of $p_\bot$ corresponds to the decrease of $\mu$. The perpendicular compression corresponds to an increase of $v_l$.  Then it proceeds with eq. (79) which I discuss further. }


\subsection{Gain and loss of energy} 

%by non-relativistic particles}



%It is clear from Eq. (\ref{pper}) that the perpendicular acceleration is also related to the motion of the magnetic mirrors at the scale $l_{\|}$. One can also notice that shrinking distance $l_{\|}$ one can increase the parallel momentum $p_{\|}$ of the CR interacting with the magnetic bottle, while at the same time, the perpendicular momentum $p_{\bot}$ is decreasing. 

%Consider the change in terms of squared perpendicular and parallel momenta of the particle. 
When the magnetic field is compressed in the perpendicular direction, according to Eq. (\ref{pper}), 
%the change of the squared perpendicular momentum is
we have 
\begin{equation}\label{eq:persqr}
\frac{ \dot{p^2_\bot}}{p^2_\bot}=\frac{2\dot{p_{\bot}}}{p_{\bot}}=-\frac{v_{l,\|}}{l_\|},
\end{equation}
where $\dot{...}$ denotes the time derivative.
Compression of
%in the direction perpendicular to 
the magnetic field entails increase of the squared perpendicular momentum,
which corresponds to energy gain. 
In an incompressible medium, compression in the perpendicular direction {with respect to the magnetic field} is inevitably accompanied by the expansion in the parallel direction. 
%and vice versa. 
In the parallel direction,  the longitudinal invariant of a mirroring particle is 
\begin{equation}
J_2=\int^b_a p_{\|} d l_\|,
\end{equation}
where  
%the component of particle momentum parallel to magnetic field
$a$ and $b$ are the two turning points.
%is not changed 
%(see Anosov \& Florinskii 1988).  
Therefore, the expansion in the parallel direction results in the loss of energy, with
%of the parallel component due to the preservation of the adiabatic invariant:
\begin{equation}
\frac{\dot{p^2_\|}}{p^2_\|}=\frac{2\dot{p_{\|}}}{p_{\|}}=2\frac{v_{l,\|}}{l_\|}.
\label{E2}
\end{equation}
%{\xu where $v_{l,\|} = - \dot{l_\|}$.}
By contrast, when the magnetic field is expanded in the perpendicular direction, 
energy loss occurs in the perpendicular direction, while energy gain occurs in the parallel direction. 

 
 
 Given Eq. (\ref{eq:persqr}) and Eq.(\ref{E2}), the change of the squared total momentum of a {particle}, i.e.,  $p^2=p_\|^2+p_\bot^2$, {under the compression and expansion in different directions} is 
 \begin{equation}
 \dot{p^2}=-(p^2_\bot-2p^2_\|)\frac{v_{l,\|}}{l_{\|}}.
 \label{Etot}
 \end{equation}
As $p_\|=p \mu$ and $p_\bot=p \sqrt{1-\mu^2}$, Eq. (\ref{Etot}) can be rewritten as 
\begin{equation}
\frac{d p^2}{dt}=-p^2(1-3\mu^2) \left(\frac{v_{l,\|}}{l_\|}\right).
\label{p2}
\end{equation}
It shows that when encountering an eddy, 
the total energy gain or loss of a particle depends not only on the perpendicular compression/expansion of the magnetic field, but also on 
the sign of 
$(1-3\mu^2)$. 




We rewrite {Eq. (\ref{p2})} to have the evolution of the total CR momentum,
\begin{equation}
\frac{d p}{dt}= -p \frac{1-3\mu^2}{2} \left(\frac{v_{l,\|}}{l_\|}\right).
\label{p1}
\end{equation}
{Formally}, the above equation is similar to Eq. (\ref{pper}) for $p_\perp$ apart from the %fact that in Eq. (\ref{p1}) $\mu$ is a variable that depends on $p_\|$ to $p_\bot$ ratio. 
{additional term $(1-3\mu^2)$.}
%It is also similar to the expression for acceleration in a highly compressible medium 
%(see \citealt{CL06}). 
%Eq. (\ref{pper}) differs from the latter expressions as it represents only part of the process rather than the change of the total $p$. On the contrary, 
%Eq. (\ref{p1}) represents the evolution of the total $p$. Thus, 
With $\mu<1/\sqrt{3} \approx 0.58$, 
the acceleration is dominated by that in the perpendicular direction.
The stochastic interaction of a particle with eddies with magnetic compression or expansion
%the corresponding {stochastic} change of the {sign of} $v_{l\|}/l_\|$ 
induces the {stochastic} increase of $p$.
%similar to {the stochastic acceleration} in a highly compressible medium \citep{CL06}.
%However, 
%in our case compression of gas is not necessary for the acceleration.
The acceleration tends to cause the decrease of $\mu$ and thus can be self-sustained with the mirroring condition always satisfied. 
%However, our setting has no compressions, and the bouncing diffusion acts instead of one induced by gyroresonance. 
%The angular momentum increase corresponds to the requirement $\mu<1/\sqrt{3}\approx 0.58$. In the opposite case of 
With $\mu>1/\sqrt{3}$, 
the acceleration is attributed to the stochastic increase of 
$p_\|$.
With the increase of $\mu$, when $\mu>\mu_c$,
particles are not subject to mirroring, and the inefficient scattering alone cannot confine particles for them to be efficiently accelerated. 
Therefore, 
to ensure the mirror diffusion and mirror acceleration, 
in the following analysis, we consider 
the upper limit of $\mu$ as
\begin{equation}
    \mu_{ca}=min\Big[\mu_{c}, \frac{1}{\sqrt{3}}\Big],
    \label{MUC}
\end{equation}
where $\mu_c$ is given by Eq. \eqref{mualeph}.


We note that Eq. \eqref{p1} would lead to 
\begin{equation}
   \frac{dp}{dt} = 0 , 
\end{equation}
when it is averaged over an isotropic distribution with $\langle \mu^2 \rangle = 1/3$. 
As expected, efficient scattering cannot lead to acceleration
(see Section \ref{ssec:scatvsmirror}).
As scattering diffusion was adopted in 
\citet{CL06},
in their model 
significant gas compression is needed 
for acceleration to happen. 







%\subsection{Evolution of cosine of pitch angle}
%\label{sec:muchange}


%{\bf AL Consider the compression of magnetic field (whatever $v_{l}/l$ is the convention adopted.) In this situation physically $\mu^2$ should decrease, i.e. the negative derivative. This should correspond to the above equation. This is fine... No problem. All my other staff about attractors is just junk.}



%\begin{equation}
%    \frac{d\mu^2}{dt}\approx (\mu^2-\mu^2_1)(\mu^2-\mu^2_2) \frac{v_{l,\|}}{l_\|},
%    \label{dmu^2}
%\end{equation}
%where
%\begin{equation}
%    \mu_1^2=\frac{7+\sqrt{13}}{12},
%\end{equation}
%and 
%\begin{equation}
%    \mu_2^2=\frac{7-\sqrt{13}}{12}.
%\end{equation}
%which are two critical points of Eq. (\ref{dmu^2}). The attractor point corresponds to 
%\begin{equation}
%    \mu_1 \approx 0.53,
%\end{equation}
%while the repellor point is
%%\begin{equation}
%    \mu_2 \approx 0.94.
%\end{equation}

%Consider the perpendicular compression of the magnetic field, i.e.,  $v_{l,\|}/l>0$. In the interval of small $\mu$, i.e., $[0,\mu_1]$, according to Eq. (\ref{dmu^2}), the derivative of $\mu^2$ is positive, inducing the increase of $\mu$. This tends to move the particles from this domain of small $\mu$ to the value of $\mu_1$. According to Eq. (\ref{p1}) this decreases the efficiency of particle acceleration to $\approx 0.16$ of its maximal value corresponding to $\mu=0$.

%For the interval $[\mu_2, 1]$ the derivative of $\mu^2$ (see Eq. (\ref{dmu^2}) is positive, which means that $\mu$ increases towards its maximal value, i.e., $\mu \rightarrow 1$. This is the region of the stable decrease of $p_\bot$. However, it corresponds to a very limited range of $\mu$. 

%For the interval $[\mu_1, \mu_2]$, the derivative of $\mu_2$ is negative and the value of $\mu$ decreases. This tends to bring the particles into the attractor point $\mu_1$. Over this interval, there is a region $[1/\sqrt{3}, \mu_2]$ where the particle momentum decreases, according to Eq. (\ref{p1}), but according to Eq. (\ref{dmu^2}) $\mu$ grows. This interval of $\mu$ presents a peculiar feature of the dynamics of $\mu$ as the perpendicular component of angular momentum $p_\bot$ decreases. Still, the corresponding decrease of $\mu$ tends to bring $p$ into the regime of increasing $p_\bot$. As the increase of $p_\bot$ corresponds to the betatron acceleration, we term this effect {\it betatron acceleration instability}. Due to this instability, the predominance of particles with larger $p_\bot$ is expected. {\bf must be tested whether this asymmetry exists in acceleration.}

%Consider magnetic field perpendicular expansion on the scale $l$, i.e., $v_{l,\|}/l<0$. The dynamics of $\mu^2$ (see Eq. (\ref{dmu^2}) is reversed. The attractor point, in this case, is $\mu_2$. The particles from the domain $[\mu_1,\mu_2]$ and from the domain $[\mu_2, 1]$ are pushed towards $\mu_2$. The particles from the interval $[\mu_1, 1]$ move towards $\mu=1$. For $\mu<1/\sqrt{3}$, the particles decrease their energy of perpendicular motion while increasing the energy corresponding to their parallel motion. 

%We have discussed earlier that the compressions and expansions proceed in a random walk fashion for the bouncing acceleration. The dynamics of $\mu$ for $\mu<1/\sqrt{3}$ ensures that the subset of particles is kept by bouncing acceleration in the bouncing diffusion regime. This is definitely true for particles with $\mu<\mu_1$. Such particles cannot leave the regime of bouncing acceleration without the pitch angle of TTD scattering. The dynamics of particles in the domain $[\mu_1, \mu_2]$ are more complicated and require further studies. We also expect the bouncing acceleration to be valid in this range. The particles with $\mu$ in the range $[\mu_2, 1]$ are bound to experience the acceleration in the parallel direction that invalidates the bouncing acceleration. 








\subsection{Fast and slow parallel diffusion}

\subsubsection{Relevant timescales}

Next we will use Eq. (\ref{p1}) to evaluate the stochastic acceleration induced by {mirror} diffusion.
%{\xu During the mirror diffusion, particles stochastically interact with slow-mode eddies with compression or expansion of magnetic fields.}
%The necessary requirement for acceleration is the random walk in the force that is experienced by the energetic particle. This is achieved as the bouncing particles sample different slow-mode eddies as they are reflected back and forth by magnetic mirrors. 

%The increment of the perpendicular momentum depends on the way that the CR interacts with the electric field of eddies. As was shown in LX21 CRs experience bouncing diffusion which for incompressible media can be many orders more slower than the diffusion arising from resonance scattering. 

For a CR to {cross} an eddy of parallel size $l_\|$, the required time is
\begin{equation}
    \Delta t_d (\mu) \approx \frac{l_\|^2}{D_\| (\mu)},
    \label{td}
\end{equation}
where $D_\|(\mu)$ is $\mu$-dependent parallel diffusion coefficient of mirror diffusion given by Eqs. \eqref{Dpar} and \eqref{Dpar2}.
To define {different regimes of particle acceleration}, $\Delta t_d(\mu)$ should be compared with the {lifetime} of the slow mode eddy:
\begin{equation}
    \Delta t_e \approx \frac{l_\|}{V_A}.
    \label{te}
\end{equation}
Therefore, we can define 
%As a result two regimes of acceleration are feasible, namely, 
the regime of {\it slow diffusion} with $\Delta t_d (\mu) > \Delta t_e$ 
and {\it fast diffusion} with $\Delta t_d (\mu)< \Delta t_e$. 
Apparently, both $\Delta t_e$ and $\Delta t_d (\mu)$ depend on the scale of turbulent eddies.
{Thus whether the diffusion is slow or fast 
depends on the range of length scales of interest.} 



{The momentum diffusion results from the stochastic interaction of particles with eddies with magnetic compression/expansion.} To find the momentum diffusion coefficient
\begin{equation}
    D_p=\frac{(\Delta p)^2}{\Delta t},
    \label{D0}
\end{equation}
one needs to know the {change} of particle momentum $\Delta p$ {after interaction with an eddy of parallel size $l_\|$,}
\begin{equation}
    \Delta p\approx \frac{dp}{dt}\Delta t.
    \label{deltap}
\end{equation}
By substituting Eq. (\ref{p1}), it is easy to obtain
\begin{equation}
    \Delta p \approx - \frac{p_{eff} v_{l,\|}}{l_\|}\Delta t 
    \label{deltap1}
\end{equation}
where
\begin{equation}
    p_{eff}= p\frac{1-3\mu^2}{2} >0.
    \label{peff}
\end{equation}
We next derive $D_p$ in different diffusion regimes. 
%The increment of momentum $\Delta p$ depends on the time of the nature of the particle-turbulence interaction, and we distinguish {\it fast parallel diffusion} and {\it slow parallel diffusion} limiting cases.
%{\xu Here} we consider the range 
%\begin{equation}
%    \mu<\frac{1}{\sqrt{3}},
%\end{equation}
%to ensure $p_{eff}>0$. 
%For $1/\sqrt{3}<\mu<\mu_c$, the parallel component of the momentum increases, increasing $\mu$ and making bouncing acceleration impossible. 
%{\xu To allow for mirror diffusion and mirror acceleration,} 
%the upper limit of $\mu$ is
%\begin{equation}
%    \mu_{ca}=min\Big[\mu_{c}, \frac{1}{\sqrt{3}}\Big],
%    \label{MUC}
%\end{equation}
%where $\mu_c$ is given by Eq. \eqref{mualeph}.
%For $\aleph^{1/2}>1/\sqrt{3}$, we observe that particles over a significant range of $\mu$ experience bouncing acceleration.



\subsubsection{Fast parallel diffusion} 

The fast parallel diffusion (FPD) of particles over a scale $l_\|$ corresponds to the requirement:
\begin{equation}
    \Delta t_d (\mu) < \Delta t_e,
    \label{less0}
\end{equation}
which is equivalent to 
(Eqs. (\ref{td}) and (\ref{te})) 
\begin{equation}
    D_\| (\mu) > l_\| V_A.
    \label{geq}
\end{equation}
 %This means that the eddy contracts on the time scale $\Delta t_e$, but the particle leaves the eddy on a shorter time scale $\Delta t_d$. The latter time is 
 As the time over which particles sample the magnetic compression/expansion within the eddy, 
 $\Delta t_d(\mu)$ enters
 Eq. (\ref{D0}), yielding 
\begin{equation}
    D_p^{FPD} (\mu) \approx \left(\frac{dp}{dt}\right)^2 \Delta t_d (\mu).
    \label{D0f}
\end{equation}
By further using Eqs. \eqref{td} and \eqref{deltap1}, 
we have 
\begin{equation}
    D_p^{FPD} (\mu)\approx \frac{v_{l,\|}^2 p_{eff}^2}{l_{\|}^2}\frac{l_{\|}^2}{D_\|(\mu)}=\frac{v_{l,\|}^2 p_{eff}^2}{ D_\| (\mu)},
    \label{Dpmu}
\end{equation}
where $D_\| (\mu)$ is given by Eqs. \eqref{Dpar} and \eqref{Dpar2}.




\subsubsection{Slow parallel diffusion}

 The case of {slow parallel diffusion (SPD)} is different.  For particles with
\begin{equation}
    \Delta t_d > \Delta t_e,
    \label{ne}
\end{equation}
%the diffusion over scale $l_\|$ is longer than the evolution of the slow mode on the scale $l_\|$, 
i.e.,
\begin{equation}
    D_\| (\mu) < l_\| V_A,
    \label{leq}
\end{equation}
%For slow parallel diffusion of particles, the momentum increment is given by Eq (\ref{deltap}), where one should use $\Delta t_e$ for the time increment. The time increment that enters the momentum diffusion in 
{$\Delta t_d (\mu)$ still enters $D_p$
in Eq. (\ref{D0}),
but $\Delta t_e$ enters $\Delta p$ in Eq. \eqref{deltap}.}
Thus we have 
\begin{equation}
\begin{aligned}
    D_p^{SPD} (\mu) &\approx \left(\frac{dp}{dt}\right)^2 \frac{(\Delta t_e)^2}{\Delta t_d (\mu)}  \\
  &  =\left[ \left(\frac{dp}{dt}\right)^2 \Delta t_d (\mu)\right]\left(\frac{\Delta t_e}{\Delta t_d (\mu)}\right)^2\\
  &  = D_p^{FPD}(\mu) \left(\frac{\Delta t_e}{\Delta t_d (\mu)}\right)^2.
    \label{Dslow0}
\end{aligned}
\end{equation}
%where the expression in square brackets coincides with the expression for the case of the fast diffusion in Eq. (\ref{D0}). 
As the ratio $\Delta t_e / \Delta t_d(\mu) <1$,
there is 
\begin{equation}
    D_p^{SPD} (\mu) < D_p^{FPD} (\mu).
\end{equation}
%of times in round brackets is $\leq 1$ which decreases the momentum diffusion in the slow parallel diffusion case. 
%Substituting in Eq. (\ref{Dslow0}) the expressions for $\Delta t_e$ and $\Delta t_d$ (see Eq. (\ref{te}) and (\ref{td})), one gets
By using Eqs. \eqref{td}, \eqref{te}, and \eqref{Dpmu}, we have 
\begin{equation}
    D_p^{SPD} (\mu) \approx \left[ \frac{v_{l,\|}^2 p_{eff}^2}{ D_\| (\mu)}\right]\left(\frac{D_\|^2 (\mu)}{V_A^2 l_\|^2}\right).
    \label{Dslow1}
\end{equation}
%where the relationship between the parallel diffusion and the evolution of the slow mode is explicitly written.
According to Eq. (\ref{Dpar}), $D_{\|} (\mu)$ strongly depends on $\mu$. 
Thus we expect that the scales over which the parallel diffusion is fast or slow 
depend on $\mu$. 
%This is the point that we study below. 




\subsection{Acceleration for fast parallel diffusion}


%The dependence of $p_{eff}$ on $\mu$ for the range of $\mu<\mu_c$ is relatively small and is not significant for our analysis. 

%Eq. (\ref{Dpmu}) presents the momentum diffusion for a particle with a given $\mu$ that interacts with slow mode turbulent motions at the given scale $l_{\|}$. The regime of fast diffusion starts when   
%  \begin{equation}
%    \Delta t_d = \Delta t_e,
 %   \label{tequal}
%\end{equation}
%or
%\begin{equation}
%    D_\| (\mu) \approx l_\| V_A.
%    \label{balance}
%\end{equation}
%and is valid for parallel scales smaller than this. 

For a given $\mu$, 
the FPD is applicable for scales smaller than a critical parallel scale $l_{\|,c}(\mu)$ corresponding to $\Delta t_d(\mu) = \Delta t_e$, 
which can be obtained by using Eqs. \eqref{Dpar}, \eqref{Dpar2}, 
\eqref{td}, and \eqref{te},
%by substituting the expression for $D_\| (\mu)$ given by Eq. (\ref{Dpar}) in Eq. (\ref{balance}):
\begin{equation}
    l_{\|,c} (\mu)= \frac{D_{\|} (\mu)}{V_A} \approx L  \aleph_s^{-2} \left(\frac{v_p}{V_A}\right) \mu^{5}
    \label{lb}
\end{equation}
at $\mu_{min}<\mu<\mu_{ca}$, and 
\begin{equation}
   l_{\|,c}(\mu) \approx   l_{\mu,min} \left(\frac{v_p}{V_A}\right)\mu
\end{equation}
at $\mu<\mu_{min}$.
{For $l_{\|,c}(\mu)$ to be larger than $l_{\|,d}$, 
the above equations constrain 
\begin{equation}
   \mu > \mu_{min,FPD,1} = \left(\frac{l_{\|,d}}{L}\frac{V_A}{v_p} \aleph_s^{2}  \right)^\frac{1}{5}
\end{equation}
at $\mu_{min} < \mu < \mu_{ca}$, and 
\begin{equation}
   \mu> \mu_{min,FPD,2} = \frac{l_{\|,d}}{l_{\mu,min}}\frac{V_A}{v_p}
\end{equation}
at $\mu< \mu_{min}$.}
We see that when 
\begin{equation}
     v_p > \Big(\frac{l_{\|,d}}{l_{\mu,min}}\Big)^\frac{5}{4} \Big(\frac{l_{\|,d}}{L}\Big)^{-\frac{1}{4}}
     \aleph_s^{-\frac{1}{2}} V_A,
\end{equation}
there is $\mu_{min,FPD,1}>\mu_{min,FPD,2}$ and vice versa.
The minimum $\mu$ for the acceleration in the FPD regime is 
\begin{equation}
  \mu_{min,FPD} = \max[\mu_{min,FPD,1},\mu_{min,FPD,2}].
\end{equation}


For $l_{\|,d}<l_\|<l_{\|,c} (\mu)$, 
the momentum diffusion coefficient at a given $\mu$ is 
(Eqs. \eqref{equns} and \eqref{Dpmu})
    \begin{equation}
    D_p^{FPD} (\mu) 
    %\approx \frac{v_{l\|}^2 p_{eff}^2}{D_{\|}(\mu)}
    \approx \frac{V_A^2 p_{eff}^2 l_\|}{ D_\| (\mu)L} \aleph_s^2,
    \label{Dmufast}
\end{equation}
where the scaling of the slow modes (see Eq.(\ref{bslow})) 
\begin{equation}
    v_{l,\|}\approx V_{L,slow} \left(\frac{l_\|}{L}\right)^{1/2}
    \label{scalslow}
\end{equation}
is used.
%If the turbulence parallel dissipation scale is $l_{\|, min}$, this introduces the minimal value of the pitch angle for which fast diffusion regime is feasible. The parallel diffusion depends on whether the diffusion coefficient is given by Eq. (\ref{Dpar}) or Eq. (\ref{Dpar2}).  In the first case, combining Eq. (\ref{Dpar}) and Eq. (\ref{balance}), it is possible to get
%\begin{equation}
%    \mu_{min1}\approx \left(\frac{l_{\|,min}}{L}\right)^{1/5} \aleph_s^{2/5} \left(\frac{V_A}{v_p}\right)^{1/5}.
 %   \label{mmin}
%\end{equation}
%The combination of Eq. (\ref{Dpar2}) and Eq. (\ref{balance}) provides:
%\begin{equation}
%    \mu_{min2}\approx \frac{V_A}{v_p}.
%    \label{mumin2}
%\end{equation}
%The boderline between two regimes of diffusion depends on the particle velocity $v_p$. For
%\begin{equation}
%    v_p<V_A \left(\frac{L}{l_{\|, min}}\right)^{1/4} \aleph_s^{-1/2},
%\end{equation}
%$\mu_{min2}>\mu_{min1}$ and vice versa.
%For $\mu$ less than
%\begin{equation}
%    \mu_{min}=max[\mu_{min1}, \mu_{min2}]
%\end{equation}
%the parallel diffusion is slower than the period of the slow mode. 
%Thus the regime of FPD is valid for $\mu \in [\mu_{min}, \mu_c]$ where $\mu_c$ is given by Eq. (\ref{mualeph}).
The resulting  averaged momentum diffusion coefficient is 
\begin{equation}
\begin{aligned}
    \bar{D}_p^{FPD} (\mu)
    &= \int_{l_{\|, d}}^{l_{\|, c} (\mu)} D_p^{FPD} (\mu) \frac{dl_\|}{l_{\|}} \\
  & \approx \frac{V_A^2 p_{eff}^2}{D_\| (\mu) L} \aleph_s^2 \left[l_{\|,c}(\mu)-l_{\|,d}\right].
    \label{Dbarf}
\end{aligned}
\end{equation}
In the case with $\mu\gg \mu_{min,FPD}$ and the scale $l_{\|,c}(\mu)\gg l_{\|, d}$, we approximately have
\begin{equation}
    \bar{D}_p^{FPD} (\mu)\approx \frac{V_A p_{eff}^2}{ L} \aleph_s^2, 
    \label{Dbarf}
\end{equation}
where $l_{\|,c}(\mu) = D_\|(\mu)/V_A$ is used. 
{Its dependence on $\mu$ only comes from the dependence of $p_{eff}$ on $\mu$.}
We see that among the eddies of different parallel sizes, the ones with
the lifetime comparable to the particle diffusion time dominate the acceleration.



%A particular feature of Eq. (\ref{Dbarf}) is that the dependence on $\mu$ disappeared. This is a consequence of the fact that $l_{\|, b} (\mu)$ is proportional to $D_\| (\mu)$. The fast diffusion takes place for 
%$\mu \in [\mu_{min}, \mu_c]$. 




\subsection{Acceleration for slow parallel diffusion with $\mu \in (\mu_{min,FPD}, \mu_{ca})$}

Within the range $\mu \in (\mu_{min,FPD}, \mu_{ca})$, 
{the parallel diffusion over scales $l_\|>l_{\|,c} (\mu)$ falls in the SPD regime. 
The averaged momentum diffusion coefficient in the SPD regime is }
\begin{equation}
    \bar{D}_p^{SPD} (\mu)\approx \int_{l_{\|, c} (\mu)}^{L} D_p^{SPD} (\mu) \frac{dl_\|}{l_{\|}}.
\end{equation}
%*************************
%where the upper limit of integration 
%\begin{equation}
%    L_{\star}=L\aleph^{-2}\mu_c^4,
%\end{equation}
%is obtained by using the scaling of the slow modes given by Eq.(\ref{bslow}) was used together with Eq. (\ref{mubk}).
%***************************
By using Eqs. (\ref{Dslow1}) and \eqref{scalslow}, we have 
\begin{equation}
    \bar{D}_p^{SPD} (\mu)\approx \frac{D_\| (\mu) p_{eff}^2}{L} \aleph_s^2 \left[\frac{1}{l_{\|, c}}-\frac{1}{L}\right].
\end{equation}
In the case with 
$l_{\| c}\ll L$, by using $l_{\|,c}(\mu) = D_\|(\mu)/V_A$, one can further write
\begin{equation}
    \bar{D}_p^{SPD} (\mu)\approx \frac{V_A p_{eff}^2}{ L} \aleph_s^2,
     \label{Dbars}
\end{equation}
which is the same as $\bar{D}_p^{FPD}(\mu)$ in Eq. \eqref{Dbarf}.
It shows that 
in our approximate treatment, { 
despite the opposite dependence of $D_p^{FPD}(\mu)$ (Eq. \eqref{Dpmu}) and $D_p^{SPD}(\mu)$ (Eq. \eqref{Dslow1}) on $D_\|(\mu)$, 
both $\bar{D}_p^{FPD}(\mu)$ and $\bar{D}_p^{SPD}(\mu)$ are dominated by the stochastic acceleration at $l_{\|,c}(\mu)$
and do not depend on $D_\|(\mu)$. 
The 
fast and slow parallel diffusion equally contribute to the momentum diffusion and acceleration.} 
%For $\mu \in [\mu_{min}, \mu_c]$, these contributions should be added to get the total diffusion coefficient for the scales in this range. 

%One can note that the expressions for the acceleration for the slow and fast regime, i.e., Eq. (\ref{Dslow1}) and Eq. (\ref{Dpmu}), have the opposite dependence on the parallel diffusion coefficient $D_{\|} (\mu)$. Nevertheless, the contribution in both cases is similar and dominated by scales in the vicinity of the border parallel scale $l_{\|,b}$. At the same time, the acceleration does not depend directly on the value of $l_{\|,b}$ given by Eq. (\ref{lb}). The only required condition is that for every $\mu$, there exists $l_{\|,b}$ that satisfies Eq. (\ref{balance}). 


\subsection{Acceleration for slow parallel diffusion with $\mu \in (0, \mu_{min,FPD})$}

For  $\mu \in (0, \mu_{min,FPD})$, the diffusion over all length scales $[l_{\|, d}, L]$ is always in the SPD regime, with 
\begin{equation}
  D_\| (\mu) < l_{\|,d} V_A=D_\|(\mu_{min,FPD}) .
  \label{eq:mumumin}
\end{equation}
 %Slow mode fluctuations on the scales $[l_{\|, min}, L]$ interact with the particle and for all the scales, inequality given by Eq. (\ref{ne}) is satisfied. 
{Given $l_{\|,d} \ll L$,} we have 
\begin{equation}
\begin{aligned}
    \bar{D}_p^{SPD} (\mu) &= \int_{l_{\|, d}}^{L} D_p^{SPD} (\mu) \frac{dl_\|}{l_{\|}}\\
    &\approx \frac{D_\| (\mu) p_{eff}^2}{L l_{\|,d}} \aleph_s^2\\
    &= \frac{V_A p_{eff}^2}{L} \aleph_s^2 \frac{D_\| (\mu)}{D_\| (\mu_{min,FPD})},
    \label{eq:Dmin}
\end{aligned}
\end{equation}
{which depends on $D_{\|}(\mu)$.}
As $D_\| (\mu) < D_\| (\mu_{min,FPD})$,
{$\bar{D}_p^{SPD}(\mu)$ with $\mu<\mu_{min,FPD}$ is smaller than that at a larger $\mu$ (Eq. \eqref{Dbars}).}
It indicates that 
the acceleration of particles with $\mu<\mu_{min,FPD}$ is less efficient than that with $\mu>\mu_{min,FPD}$. 



%************************
%Various damping processes can affect $l_{\|,d}$. However, our results show that this can only marginally affect the acceleration process. As the bouncing acceleration decreases $\mu$, more particles can be concentrated in the domain $\mu<\mu_{min}$. Our calculations show that will decrease the efficiency of the acceleration, but will not shut down the bouncing acceleration for such particles.
%************************


\subsection{Total momentum diffusion coefficient}


{The averaged total momentum diffusion coefficient $\bar{D}_p^{total}(\mu)$ is given by 
$\bar{D}_p^{SPD}(\mu)$ in Eq. \eqref{eq:Dmin} at $\mu< \mu_{min,FPD}$.}
For $\mu \in (\mu_{min,FPD}, \mu_{ca})$, 
the SPD and FPD processes should be both taken into account, leading to 
\begin{equation}\label{eq:totpu}
    \bar{D}_p^{total} (\mu)\approx D_p^{FPD} (\mu) + D_p^{SPD} (\mu)\approx \frac{2V_A p_{eff}^2}{ L} \aleph_s^2,
\end{equation}
where we sum up the results in Eq. (\ref{Dbarf}) and Eq. (\ref{Dbars}).

%As the acceleration happens for the range of $\mu$ the diffusion in terms of total momentum for an ensemble of particles can be obtained by averaging over the distribution of pitch angle 
{Furthermore, by averaging over pitch angles, we have} 
\begin{equation}
\begin{aligned}
    \bar{D}_p^{total}
    &=\frac{1}{\int_0^{\pi/2}f d\theta}
    \int_{\theta_{ca}}^{\pi/2} \bar{D}^{total}_p (\mu(\theta)) f (\mu(\theta)) d\theta \\
    &=\frac{1}{\int_0^{\pi/2}f d\theta} \int_0^{\mu_{ca}} \bar{D}_p^{total} (\mu) f (\mu) \frac{d\mu}{\sqrt{1-\mu^2}}.
    \label{Dpfast}
\end{aligned}
\end{equation}
%For typical settings $\mu_{min}\ll 1$ and therefore the contribution of particles with $\mu<\mu_{min}$ to the total momentum diffusion:
%\begin{equation}
%\begin{aligned}
%     \delta \bar{D}_p^{total}
%     &\approx \frac{2}{\pi} \int^{\mu_{min}}_{0} D_p^{total} (\mu) f (\mu) \frac{d\mu}{\sqrt{1-\mu^2}}\\
%     &\sim \mu_{min} D_p^{total} (\mu_{min}) f (\mu_{min})
%\end{aligned}
%\end{equation}
%is negligible compared to the contribution from $\mu>\mu_{min}$.
%It was discussed in LX21 that a special case of a uniform distribution of pitch angles can be present in particular settings, e.g. near the sources of energetic particles, in the presence of resonance scattering by fast modes. 
%If the distribution of particles is uniform in terms of $\mu$, the total diffusion momentum diffusion coefficient can be estimated as
For simplicity, in the case with a uniform pitch angle distribution, i.e., $f=1$,
\footnote{Note that the uniform pitch angle distribution considered here is different from the fast pitch-angle isotropization achieved by efficient pitch-angle scattering
during the interaction of a particle with an eddy discussed in Section \ref{ssec:scatvsmirror}.}
there is 
(Eqs. \eqref{peff} and \eqref{eq:totpu})
\begin{equation}\label{eq:dptamca}
\begin{aligned}
    \bar{D}_p^{total}
    &\approx \frac{1}{\pi} \frac{V_A p^2}{L} \aleph_s^2 \int_{\mu_{min,FPD}}^{\mu_{ca}} (1-3\mu^2)^2 \frac{d\mu}{\sqrt{1-\mu^2}} \\
    &\approx \frac{1}{\pi} \frac{V_A p^2}{L} \aleph_s^2 \int_{\mu_{min,FPD}}^{\mu_{ca}} (1-3\mu^2)^2 d\mu \\
    & \approx  \frac{1}{\pi} \frac{V_A p^2}{L} \aleph_s^2 \left(\frac{9\mu_{ca}^5}{5}-2\mu_{ca}^3+\mu_{ca}\right)\\
    & \approx \frac{1}{\pi} \frac{V_A p^2}{L} \aleph_s^2 \mu_{ca},
\end{aligned}    
\end{equation}
{where we assume that $\mu_{min,FPD}$ is close to $0$ and consider $\mu_{ca}\leq 1/\sqrt{3}$.}
%The integration formally provides
%\begin{equation}
 %   D_p\approx...
%\end{equation}
The corresponding acceleration timescale can be estimated as  
\begin{equation}\label{tauacc}
\begin{aligned}
    \tau_{acc} &\approx \frac{p^2}{\bar{D}_p^{total}} \\
    &\approx \pi \frac{L}{V_A}  \aleph_s^{-2} 
    \left(\frac{9\mu_{ca}^5}{5}-2\mu_{ca}^3+\mu_{ca}\right)^{-1}\\
    & \approx \pi \frac{L}{V_A} \aleph_s^{-2} \mu_{ca}^{-1}.
 %   = \frac{L}{V_A}\left(\frac{V_A}{V_{L,slow}}\right)^2
\end{aligned}
\end{equation}
{We see that both $\bar{D}_p^{total}$ in Eq. \eqref{eq:dptamca} and 
$\tau_{acc}$ in Eq. \eqref{tauacc} are independent of $D_{\|}(\mu)$.}
This is due to the fact that 
irrespective of $D_\|(\mu)$, it is the eddies with their lifetime equal to the diffusion time corresponding to $D_\|(\mu)$ that dominate the acceleration. 





%\subsection{Effect of finite Larmor radius}

%The interaction of particles with magnetic mirrors gets modified when the effects of the finite Larmor radius of a particle become important. The process of bouncing is modified when the effect of finite Larmor radius is accounted for (XL20). For $\mu<\mu_{L,s}$, where
%\begin{equation}
%    \mu_{L,s} \approx \aleph_s^{1/2} \left(\frac{r_L}{L}\right)^{1/4},
%\end{equation} 
%the corresponding coefficient for bouncing diffusion, according to LX21 is
%\begin{equation}
%    D_{\|} (\mu)\approx v_p \mu r_L. ~~~~~~ \mu<\mu_{L, s}.
%    \label{eq:rL}
%\end{equation}

%The simplest modification of the earlier results is applicable if the value of $\mu_{L,s}$ is less than $\mu_{min}$. This means that the change in the diffusion coefficient happens at $\mu$ for the range of slow parallel diffusion. Then the diffusion coefficient in Eq. (\ref{eq:Dmin}) up to $\mu_{min,s}$ is given by Eq. (\ref{eq:rL}) and for
%$\mu_{min,s}<\mu<\mu_{min}$ is given by Eq. (\ref{Dpar}). 

%If $\mu_{L,s}>\mu_{min}$, then the transition between the regimes of slow and fast diffusion takes place in the regime when the diffusion coefficient is given by Eq. (\ref{eq:rL}). Then 
%the value 
%\begin{equation}
%    \mu_{minL}\approx \frac{l_{\|,min}}{r_L}\frac{V_A}{v_p},
%    \label{muminL}
%\end{equation}
%should be employed instead of $\mu_{min}$ in our earlier expressions of momentum diffusion. Comparing Eq. (\ref{muminL}) and Eq. (\ref{mumin2}) one can see that the necessary condition for the finite Larmor radius to be important is that the particle Larmor radius is less than $l_{\|, min}$.

%In most cases, the contributions to total acceleration arising from $\mu<\mu_{minL}$ are marginal and can be neglected. 




\subsection{Pitch angle diffusion due to mirror acceleration}

%Mirror diffusion via interactions with magnetic mirrors preserves {\xu the first adiabatic invariant} given by Eq. (\ref{J1})
%{\xu and thus does not cause diffusion in pitch angle when $p$ is constant.}
%Thus, unless the magnetic field $B$ changes, the value of $\mu$ stays the same if the particle interacts with not moving magnetic mirrors. 
%{\xu The magnetic mirrors created in MHD turbulence are not static. Their motions 
%induce the mirror acceleration process.} 
%As shown in \S \ref{sec:muchange}, 
The mirror acceleration results in {change of $p$ and thus change of $\mu$.} 
Similar to the transit time damping (TTD), the change of $\mu$ is associated with acceleration.
But unlike the resonant TTD that causes the stochastic increase of $\mu$, 
the non-resonant mirror acceleration causes the stochastic decrease of $\mu$.





The evolution of $\mu$ can be determined by using Eq. (\ref{pper}),
\begin{equation}
    \frac{d}{dt} (p\sqrt{1-\mu^2}) \approx -\frac{1}{2}p\sqrt{1-\mu^2} \frac{v_{l,\|}}{l_\|},
\end{equation}
which provides
\begin{equation}
    \sqrt{1-\mu^2} \frac{dp}{dt} - p \frac{\mu}{\sqrt{1-\mu^2}} \frac{d\mu}{dt}\approx -\frac{1}{2}p\sqrt{1-\mu^2} \frac{v_{l,\|}}{l_\|}.
    \label{eq:deriv}
\end{equation}
The substitution of Eq. (\ref{p1}) in Eq. (\ref{eq:deriv}) provides
\begin{equation}\label{dmu2new}
    \frac{d\mu^2}{dt} =3\mu^2 (1-\mu^2)  \frac{v_{l,\|}}{l_\|}.
\end{equation}
Then the time derivative $d\mu/dt$ is 
%in bouncing acceleration 
\begin{equation}\label{eq:dmudt}
\frac{d\mu}{dt} = \frac{3}{2} \mu (1-\mu^2) \frac{v_{l,\|}}{l_\|}.
\end{equation}
%\begin{equation}
%    \frac{d\mu}{dt}\approx \frac{1}{2\mu}(\mu^2-\mu^2_1)(\mu^2-\mu^2_2) \frac{v_{l,\|}}{l}.
%    \label{dmu^21}
%\end{equation}


{Similar to the derivation of momentum diffusion coefficient, the derivation of pitch angle diffusion coefficient also depends on the diffusion regime.}
%increments of $\Delta p$ that we considered earlier, the increments $\Delta \mu$ depend on whether the diffusion is slow or fast, i.e., depends on the relation between $\Delta t_e$ and $\Delta t_d$, i.e., whether the parallel diffusion is fast or slow. 
In the FPD regime, the time interval relevant in determining $\Delta \mu$ is $\Delta t_d(\mu)$, i.e.,
$\Delta \mu \approx (d\mu/dt) \Delta t_d(\mu)$. The time step for the diffusion is also $\Delta t_d(\mu)$. Thus the pitch angle diffusion coefficient in the FPD regime is 
(Eqs. \eqref{equns}, \eqref{scalslow}, \eqref{td}, and \eqref{eq:dmudt})
\begin{equation}
\begin{aligned}
    D_{\mu\mu}^{FPD} (\mu)
    &\approx \frac{(\Delta \mu)^2 }{\Delta t} \\
    &\approx \left(\frac{d\mu}{dt}\right)^2 \Delta t_d(\mu) \\
    & = \frac{9 \mu^2(1-\mu^2)^2 V_A^2 l_\|}{4 D_\|(\mu)L}  \aleph_s^2.
    \label{Dmm}
\end{aligned}
\end{equation}
Its dependence on $\mu$ is explicitly shown.
At $\mu_{min,FPD} <\mu<\mu_{ca}$, the averaged pitch angle diffusion coefficient is 
\begin{equation}
\begin{aligned}
   \bar{D}_{\mu\mu}^{FPD}(\mu) 
   &= \int_{l_{\|,d}}^{l_{\|,c}(\mu)}D_{\mu\mu}^{FPD}(\mu) \frac{dl_\|}{l_\|} \\ 
   &\approx  \frac{9 \mu^2(1-\mu^2)^2 V_A^2}{4 D_\|(\mu)L}  \aleph_s^2 [l_{\|,c}(\mu)-l_{\|,d}].
\end{aligned}
\end{equation}
Provided $l_{\|,c}(\mu)\gg l_{\|,d}$ and using $D_{\|}(\mu) = V_A l_{\|,c}(\mu)$, we approximately have 
\begin{equation}\label{eq:appdmmf}
\bar{D}_{\mu\mu}^{FPD}(\mu) \approx 
  \frac{9 \mu^2(1-\mu^2)^2 V_A}{4 L}  \aleph_s^2 .
\end{equation}
{It does not depend on $D_\|(\mu)$.}



In the SPD regime, 
%i.e., when parallel particle diffusion is slow, the particle samples the effect of the fluid contraction for the perturbation period, 
we have $\Delta \mu\approx (d\mu/dt) \Delta t_e$. 
%The relevant time for the random walk is, as in the case of the momentum diffusion, $\Delta t_d$. 
As a result, there is 
\begin{equation}
\begin{aligned}
    D_{\mu\mu}^{SPD}(\mu)
    &\approx \left( \frac{d\mu}{dt}\right)^2 \frac{(\Delta t_e)^2}{\Delta t_d(\mu)} \\
    &= \left[\left(\frac{d\mu}{dt}\right)^2 \Delta t_d(\mu)\right] \left(\frac{\Delta t_e}{\Delta t_d(\mu)}\right)^2 \\
    & \approx D_{\mu\mu}^{FPD}(\mu) \left(\frac{\Delta t_e}{\Delta t_d(\mu)}\right)^2 \\
    & \approx 
     \frac{9 \mu^2(1-\mu^2)^2 V_A^2 l_\|}{4 D_\|(\mu)L}  \aleph_s^2 \left(\frac{D_{\|}^2(\mu)}{V_A^2 l_\|^2}\right),
\end{aligned}
\end{equation}
where Eqs. \eqref{td}, \eqref{te}, and \eqref{Dmm} are used. 
With $D_\|(\mu) < V_A l_\|$, we naturally have 
\begin{equation}
   D_{\mu\mu}^{SPD}(\mu) < D_{\mu\mu}^{FPD}(\mu).
\end{equation}
At $\mu_{min,FPD}<\mu<\mu_{ca}$, the averaged pitch angle diffusion coefficient in the SPD regime is 
\begin{equation}
\begin{aligned}
  \bar{D}_{\mu\mu}^{SPD}(\mu)
  &= \int_{l_{\|,c}(\mu)}^L D_{\mu\mu}^{SPD}(\mu) \frac{dl_\|}{l_\|} \\
  &\approx \frac{9\mu^2(1-\mu^2)^2D_\|(\mu)}{4L}\aleph_s^2\left(\frac{1}{l_{\|,c}(\mu)}-\frac{1}{L}\right).
\end{aligned}
\end{equation}
Given $l_{\|,c}(\mu)\ll L$ and using the relation $D_{\|}(\mu) = V_A l_{\|,c}(\mu)$, the above expression is approximately
\begin{equation}\label{eq:apdmmsbm}
\bar{D}_{\mu\mu}^{SPD}(\mu) \approx
\frac{9\mu^2(1-\mu^2)^2 V_A}{4L}\aleph_s^2,
\end{equation}
which is the same as 
$\bar{D}_{\mu\mu}^{FPD}(\mu)$ in Eq. \eqref{eq:appdmmf}.

At $\mu< \mu_{min,FPD}$, the averaged pitch angle diffusion coefficient is 
\begin{equation}
\begin{aligned}
  \bar{D}_{\mu\mu}^{SPD}(\mu) &=
  \int_{l_{\|,d}}^{L} D_{\mu\mu}^{SPD}(\mu) \frac{dl_\|}{l_\|} \\
&\approx   \frac{9\mu^2(1-\mu^2)^2D_\|(\mu)}{4L}\aleph_s^2\left(\frac{1}{l_{\|,d}}-\frac{1}{L}\right).
\end{aligned}
\end{equation}
Given $l_{\|,d}\ll L$, it approximately becomes 
\begin{equation}
\bar{D}_{\mu\mu}^{SPD}(\mu) \approx 
   \frac{9\mu^2(1-\mu^2)^2 V_A }{4L}\aleph_s^2\frac{D_\|(\mu)}{D_\|(\mu_{min,FPD})},
\end{equation}
where Eq. \eqref{eq:mumumin} is used. It depends on $D_{\|}(\mu)$. 
With $D_{\|}(\mu) < D_{\|}(\mu_{min,FPD})$, 
$\bar{D}_{\mu\mu}^{SPD}(\mu)$ at $\mu<\mu_{min,FPD}$ is smaller than that at a larger $\mu$ (Eq. \eqref{eq:apdmmsbm}). 
{The more efficient pitch angle diffusion at $\mu > \mu_{min,FPD}$ is caused by the more efficient acceleration. }



%It is easily observed that the time derivative of $\mu$ according to Eq. (\ref{dmu^21}) formally gets to infinity as $\mu\rightarrow 0$. This shows that in the presence of bouncing acceleration, the classical suppression of scattering at $\mu\rightarrow 0$ does not exist.

%Further treatment of the $\mu$-diffusion follows step by step the $p$-diffusion that we considered above. In the turbulent cascade, smaller eddies from $l_{\|, min}$ determined by damping to the parallel border scale $l_{\|,b}$ given by Eq. (\ref{lb}) contribute to the $\mu$-diffusion in the FPD regime. The corresponding integration is
%\begin{equation}
%    \bar{D}_{\mu\mu}^{FPD} (\mu)\approx \int_{l_{\|, min}}^{l_{\|, b} (\mu)} D_{\mu\mu}^{FDP} (\mu) \frac{dl_\|}{l_{\|}}.
%    \label{Dbarf01}
%\end{equation}

 %Similarly, the turbulent scales larger than $l_{\|,b}$ contribute to the $\mu$-diffusion in SPD regime. The integration in this case starts from $l_{\|, b}$:
 %\begin{equation}
%    \bar{D}_{\mu\mu}^{SPD} (\mu\leq \mu_{min})\approx \int_{l_{\|, min}}^{L_\star} D_{\mu\mu}^{SPD} (\mu) \frac{dl_\|}{l_{\|}}.
%   \label{eq:slowmin11}
%\end{equation}
%***************
%The difference between the $\mu$-diffusion from the $p$-diffusion is that for the latter the diffusion integrated over the distribution of $\mu$ is valuable. In the case of $\mu$-diffusion, the dependence on $\mu$ should not be integrated out. 





%\subsection{Mirror acceleration: $\mu$ dependence}

%The difference of the bouncing acceleration arises from the fact that the bouncing diffusion is a function of $\mu$.  

%due to the different energy gain/loss in parallel and perpendicular directions 

%The direct consequence of the dependence of the parallel bouncing diffusion coefficient $D_\| (\mu)$ on the pitch angle $\mu$ is that the diffusing time given by Eq. (\ref{td}) is a function of $\mu$.
%**********




\begin{comment}

\section{Mirror acceleration in a low-$\beta$ medium}

%In a low-$\beta$ medium, 
%due to the inefficient scattering by 
%Alfv\'{e}n and slow modes, 
%the pitch-angle scattering is dominated by fast modes at $\mu>\mu_c$. 
%cannot be disregarded \citep{YL04}, while the pitch-angle scattering by Alfv\'{e}n and slow modes is inefficient for $r_g\ll L$ {\xu due to the scale-dependent turbulence anisotropy}. 
%The borderline between two regimes is determined by $\mu_{c, scat}$ given by Eq. (\ref{eq:mucscat}). The acceleration of non-bouncing particles we do not consider; it is covered in CL06. We focus on bouncing acceleration for particles with $\mu<\mu_{c, scat}$. Contrary to the previous section, we consider purely compressible motions where a compression in one direction does not induce expansion in another direction. This simplifies our treatment compared to the earlier studied incompressible case.

%\subsection{General considerations}

%The acceleration physics differs in low-$\beta$. 
%In a compressible low-$\beta$ medium, 
%unlike magnetic field compressions in an incompressible medium, for fluid compressions, there is no relation similar to Eq. (\ref{p1}) between the parallel and the perpendicular momenta of particles. We employ the parallel diffusion coefficient of bouncing diffusion $D_\| (\mu)$ that, as in the earlier discussed high-beta plasma case, is a function of $\mu$. Therefore, as earlier, for a given $\mu$, we have to consider the acceleration in FPD regime that arises from particle interaction with small eddies and the acceleration in SPD regime arising from the interaction with large eddies of the MHD turbulent cascade. 



%Consider fluid compressions in magnetized media with particles performing bouncing diffusion. 
In a compressible low-$\beta$ medium,
the compression of the magnetic field in the perpendicular direction causes the increase of $p_\perp$, and the compression in the parallel direction causes the increase of $p_\|$. 
As a result, the change of the total momentum of a particle propagating through 
an eddy of size $l$ can be expressed as (see \citealt{CL06})
\begin{equation}
    \frac{dp}{dt}\approx -p \frac{{\bf \nabla} \cdot{\bf v}_l}{\beth},
    \label{eq:beth}
\end{equation}
where {\xu $v_l$ is the velocity of the compressible component of turbulence at $l$,}
%and compressible velocity with compressible component ${\bf v_l}$ 
and 
$\beth=3$ for isotropic compression. 
The minus sign in Eq. (\ref{eq:beth}) reflects the increase of $p$ due to compression and the decrease of $p$ due to expansion. 


{\xu Similar to the acceleration in an incompressible medium, for acceleration in a compressible medium, we also separately treat the acceleration process in fast and slow parallel diffusion regimes. } 

%The rest of the acceleration process is similar to that we considered earlier. Indeed, it is appropriate to distinguish the slow and fast limit of diffusion depending on whether the particle crosses the eddy faster or slower than the eddy evolves. The corresponding time scale $\triangle t_d$ is determined by the bouncing diffusion through the scale $l$ and coincides with the diffusion time scale considered for the incompressible case. The time scale of eddy evolution $\triangle t_e$ varies depending on the fluid motions we consider. Several types of compressions are considered below.

\subsection{Acceleration by slow modes}

%For the acceleration by slow modes, the effects of fluid compression are most important in the limit of low-$\beta$ plasmas. 

\subsubsection{Fast parallel diffusion}

%The condition for FPD is given by 
%Eq. (\ref{less0}).
%The direct consequence of the dependence of the parallel bouncing diffusion coefficient $D_\| (\mu)$ on the pitch angle $\mu$ is that the diffusing time given by Eq. (\ref{td}) is a function of $\mu$. 
In a low-$\beta$ medium, slow modes cause compression of the fluid. The rate of compression in an eddy of size $l$ is $v_l/l$. For simplicity, we have (Eq. \eqref{eq:beth})
\begin{equation}
    \frac{dp}{dt}\approx p \frac{v_l}{l}.
    \label{eq:pchange}
\end{equation}
%This compression is sampled over a shorter $\Delta t_d (\mu)$. This time should be used in Eq. (\ref{deltap}) to find the momentum increment over the time that a particle with pitch angle $\mu$ interacts with the eddy. As a result, one gets
%\begin{equation}
%    \Delta p \approx \frac{p v_{l}}{l}\Delta t ,
%    \label{compr:deltap1}
%\end{equation}
%the major difference of which from Eq. (\ref{deltap1}) is due to using $p$ instead of $p_{eff}$ in the latter expression. This difference arises from the fact that in the case of fluid compression, there is no compensating expansion that decreases another component of momentum. 
%The range of $\mu$ for the acceleration is $\mu<\mu_{c, scat}$, where $\mu_{c, scat}$ is given by Eq. (\ref{eq:mucscat}). Depending on the scattering rate (see Eq. (\ref{eq: appfqltu})), the range of $\mu$ for bouncing acceleration in the low-$\beta$ medium can be larger or smaller than the range of $\mu$ for bouncing acceleration in the high-$\beta$ medium that is given by Eq. (\ref{MUC}). 
The corresponding momentum diffusion coefficient in the FPD regime is
(Eqs. \eqref{td}, \eqref{D0f}, and \eqref{eq:pchange})
\begin{equation}
    D_p^{FPD}(\mu)\approx \frac{v_l^2 p^2}{D_\| (\mu)}, 
\end{equation}
which is a compressible counterpart of Eq. (\ref{Dpmu}), {\xu with $p$ used instead of $p_{eff}$.}




\subsubsection{Slow parallel diffusion}

In the SPD regime, the approach is similar to the one we applied to the incompressible medium. %The increment of momentum $\Delta p$ is obtained through the action of velocity $v_l$ at the scale $l$, see Eq. (\ref{eq:pchange}). The time of the step that determines the momentum diffusion coefficient in Eq. (\ref{D0}) is $\delta t_d$.
Thus in analogy to Eq. (\ref{Dslow1}), we
have
\begin{equation}
    D_p^{SPD} (\mu) \approx \left[ \frac{v_{l}^2 p^2}{ D_\| (\mu)}\right]\left(\frac{D_\|^2 (\mu)}{V_A^2 l^2}\right),
    \label{Dslownew}
\end{equation}




\section{SubAlfvenic and SuperAlfvenic setting}

\subsection{SubAlfvenic turbulence}



For the subAlfvenic turbulence, one can employ a more formal approach proposed in LX21, namely, the introduction of an effective injection scale $L_{eff}$:
\begin{equation}
    L_{eff}=L M_A^{-4},
\end{equation}
which does not have a physical meaning, but, if inserted instead of $L$ in the expressions, provides the right scaling for subAlfvenic turbulence statistics.

\subsection{SuperAlfvenic turbulence}

If the Larmor radius of a particle is small compared to either $l_A$ for superAlfvenic turbulence or $l_{tran}$ for subAlfvenic turbulence, then the acceleration is induced by slow modes that follow the scaling of strong turbulence.

The expected change of our earlier equations is minimal. In fact, for the superAlfvenic turbulence, instead of $L$ one should use $l_A$ given by Eq. (\ref{lA}). This is natural as $l_A$ can be viewed as the injection scale for actual MHD turbulence.

The stochastic acceleration by superAlfvenic turbulence is expected to differ for small and large scales. For scales less than the Alfven scale
\begin{equation}
    l_A=LM_A^{-3},
\end{equation}
the turbulent properties are identical to the transAlvenic turbulence with the injection scale of $l_A$. For scales, $l>l_A$, the magnetic field structure has not been carefully studied.

From general physical considerations, one can expect a weak magnetic field to follow passively fluid motions. In the subsonic setting, we expect the hydro turbulence to follow the incompressible Kolmogorov cascade. This will radically change both Alfvenic and slow mode evolution. The incompressible stretches and contractions of the magnetic field can be viewed as the extension of the Alfvenic cascade. The compressions of the magnetic field are compensated by changes in the gaseous pressure as the extension of the slow mode cascade. The hydrodynamic cascade will impose the Kolmogorov scaling on both types of magnetic perturbations. 

For magnetic field compressions, this would mean that at the injection scale of superAlfvenic turbulence, the compressible fluctuation will be the order of the total magnetic field. In terms of bouncing, that means that at the large scale, one can expect $\mu_c\approx 1$. 

The scaling of velocity fluctuations that move the magnetic mirrors is expected to be Kolmogorov: 
\begin{equation}
    v_l\approx V_L \left(\frac{l}{L}\right)^{1/3},
\end{equation}
with 



\section{Comparisons with earlier studies} 
 

\subsection{Comparison with \citet{CL06}}

A detailed study of particle acceleration by large-scale {\xu compressible} turbulent motions was performed by \citet{CL06}. 
There the acceleration by slow modes in an incompressible medium was also studied.  {\xu However, the mirror diffusion identified in LX21
was not considered. }
%Cho \& Lazarian (2006) were unaware of the bouncing diffusion process introduced only in LX21. Therefore the authors considered two alternatives. 
They found that if particles 
move ballistically along magnetic field lines, 
{\xu the acceleration due to magnetic compressions is inefficient as the particle speed is usually much larger than $V_A$.}
%and experience an electric field arising from the time variations of magnetic compressions. Due to the disparity of typical Alfven and particle velocities, this process was a very inefficient accelerator. 
{\xu When efficient pitch-angle scattering that enables slow diffusion of particles was assumed, 
%When the authors allowed efficient scattering, 
they found that the corresponding isotropic pitch angle distribution is problematic,
as the acceleration due to magnetic compression in the perpendicular direction is canceled by the deceleration in the parallel direction.}
%another problem with the acceleration mechanism. As we discussed in \S ..., the compression of the magnetic field in the perpendicular direction in incompressible limit induces the expansion of the parallel direction. As a consequence, if scattering is efficient and the distribution of particles is isotropic, the total transfer of energy from slow modes to particles vanishes. 
As a result, they concluded that the non-resonance particle acceleration by large-scale magnetic compressions takes place only due to fluid compressions {\xu in a compressible medium with efficient pitch-angle scattering.} 
However, with efficient pitch-angle scattering, the scattering-induced acceleration would become more efficient than the compression-induced acceleration.
%At the same time, the authors correctly pointed out that the necessary requirement for such acceleration is the existence of resonant pitch-angle scattering, which, in turn, is more efficient than the compressional acceleration they studied. This was contrary to the conclusions reached by earlier authors (see Ptuskin 1988, Chandran ...) and meant that the stochastic acceleration by turbulence compressions is not important unless some special conditions, e.g., shock compressions in supersonic turbulence, are involved. However, the shock compression case is special as shocks can induce more efficient first-order Fermi acceleration.

Our conclusion differs from that of \citet{CL06}. With the mirror diffusion, the energetic particles do not need to undergo efficient pitch angle scattering to enhance their interaction with slow modes. 
As a result, we find that 
%large-scale compressions' stochastic acceleration can successfully compete with pitch angle scattering. 
{\xu magnetic compressions in an incompressible medium can lead to efficient particle acceleration.}
This is especially important 
when the turbulent motions at length scales $\lesssim r_g$ are damped 
\citep{Brunetti_Laz}
%as the turbulence on small scales experiences damping, e.g., collisionless Landau damping (see Yan \& Lazarian 2003 for low-beta media and Brunetti \& Lazarian 2007 for high-beta media). 
or when the pitch-angle scattering is inefficient in an incomressible medium. 
%This damping suppresses pitch-angle scattering. On the contrary, large-scale compressions can still efficiently accelerate particles in turbulence.   






\subsection{Comparison with the acceleration arising from pitch-angle scattering}

Bouncing diffusion, unlike the the diffusion induced by pitch angel scattering does not change the adiabatic invariant. At the same time, similar to scattering, it introduces the mean free path for the particles. Compared to pitch angle scattering, bouncing diffusion rate increases with the decrease of pitch angle $\mu$ (see Eq. (\ref{Dpar})). Thus bouncing diffusion is the process that resolves the well-known paradox of earlier cosmic ray propagation theories that arises from the inefficiency of scattering for $\mu\rightarrow 0$, i.e. the so-called "90 degree scattering paradox" (Jokipii...). 

Pitch angle scattering plays a dual role in both bouncing propagation and bouncing acceleration. On one hand, scattering decreases the range of $\mu$ for which the particles are mirrored by magnetic fluctuations (Kulsrud ..., Xu \& Lazarian 2020). On the other hand, pitch angle scattering can replenish the population of particles in the range of $\mu$ corresponding to bouncing diffusion and acceleration. 

Both processes enable the stochastic acceleration arising from the large-scale turbulent compression. However, unlike the bouncing diffusion, the pitch angle scattering induces itself the stochastic acceleration that is faster than the acceleration from large-scale compressions. Indeed, the acceleration corresponding to the following momentum diffusion coefficient (Cho \& Lazarian 2006) 
\begin{equation}
    D_p^{scatt}\approx \frac{p^2 V_A}{L}\left(\frac{LV_A}{l_{mfp} v_p}\right).
\end{equation}
In terms of rough comparison, the pitch angle scattering dominates if
$LV_A>l_{mfp} v_p$, i.e. if $l_{mfp}<LV_A/v_p$. This can be true in the presence of fast modes, but is difficult to realize with slow and Alfven modes (see Yan \& Lazarian 2003). 

At the same time, pitch angle scattering shuts down the acceleration via slow modes as the fluid gets incompressible. On the contrary, the boucing diffusion, as we demonstrated in this paper, enables efficient stochastic acceleration by slow modes in incompressible fluid.  

\subsection{Comparison with acceleration by Alfven modes}
Alfven modes are incompressible; therefore, they do not induce the change of perpendicular momentum of a particle. However, the fluctuations of the length of the magnetic field line can induce the parallel momentum change due to the preservation of the second (parallel to the field) adiabatic invariant. In the presence of bouncing diffusion, the acceleration through this process can proceed for pitch angles $\mu<\mu_c$ inducing random walk increase of the particle parallel momentum.
\begin{equation}
\frac{\dot{E^{kin}_\|}}{E^{kin}_\|}=2\frac{\dot{p_{\|}}}{p_{\|}}=-2\frac{v_{\|}}{l_\|} \left(\frac{v_l}{V_A}\right)^2,
\label{E}
\end{equation}
where the small factor is present in the round brackets. 

As Alfvenic modes are incompressible, a corresponding magnetic field compression corresponds to the magnetic field line extension. Thus we expect that the momentum will evolve according to
\begin{equation}
\frac{d p}{dt}=p^2(1-3\mu^2) \left(\frac{v_{l,\|}}{l_\|}\right)\left(\frac{v_l}{V_A}\right)^2,
\label{p2A}
\end{equation}
which is different from Eq. (\ref{p1}) by the presence of the squared ratio of the wave amplitude at the scale $l$ to the Alfven velocity. 

Due to the superdiffusion of magnetic field lines, the bouncing acceleration is possible in the presence of magnetic bottles, e.g., created by the fast or slow modes, for $\mu<1/\sqrt{3}$. However, we expect this type of acceleration induced by Alfven modes to be sub-dominant to the bouncing acceleration arising from the other two types of modes. 

Alfven modes are more robust to collisionless damping (see Yan \& Lazarian 2003). However, in the absence of slow and fast modes, the magnetic field compressions are negligibly small to induce appreciable bouncing. 

The stochastic acceleration in terms of the parallel momentum is possible for Alfven modes for $\mu>1/\sqrt{3}$. If scattering limits the mean free path of particles, the stochastic acceleration is due to the scattering and not due to compressions (CL02). 


One can consider ballistic propagation, i.e., propagation without collisions.  Suppose the particle is zooming along the magnetic field without scattering. In that case, this is the Fast Parallel Diffusion (FPD) regime with an extremely large diffusion coefficient $D_\| (\mu)$. Stochastic gaining of momentum in the FPD is regulated by the momentum diffusion coefficient that can be obtained by multiplying the RHS of Eq. (\ref{Dpmu}) by $(v_l/V_A)^4$ factor, and taking into account that in the absence of bouncing $p_{eff}$ should be changed for $p$ the coefficient gets independent of $\mu$, i.e.,
\begin{equation}
    D_p^{FPD} \approx \frac{v_\|^2 p^2}{ D_\| (\mu)}\left(\frac{v_l}{V_A}\right)^4.
    \label{DpAlf}
\end{equation}
 
The contribution increases with the  coefficient goes to zero as $D_\| (\mu)\rightarrow \infty$. 

Incidentally, the latter regime corresponds to particle acceleration moving without scattering in pure Alfvenic turbulence. The process of magnetic field shrinking can be associated with turbulent reconnection in LV99. The extension of the magnetic field can be associated with local dynamo (Beresnyak ...). The balance of the two processes provides the stationary state of Alfvenic turbulence. Our considerations above show that the stochastic increase of the parallel momentum for ballistically moving particles is negligibly small for sub-Alfvenic turbulence. 



The situation is more involved in the superAlfvenic case. In the absence of collisions, the Alfvenic scale $l_A=LM_A^{-3}$ acts as the effective mean free path (Brunetti \& Lazarian 2007). This limits the value of $D_\|$ and increases the efficiency of acceleration. More importantly, in this case, one can expect the formation of extended long-lived regions of magnetic reconnection. The acceleration in such regions is expected to differ from the process discussed in this paper (see Lazarian et al. 2020). In Brunetti \& Lazarian (2016), we speculated that this could be an efficient way of cosmic ray reacceleration in galaxy clusters.

Due to this small factor, the acceleration is inefficient for sub-Alfvenic turbulence. However, the acceleration due to preserving the second adiabatic invariant can be fast in super-Alfvenic turbulence and related to stochastic reconnection acceleration in Brunetti \& Lazarian (2011) that we briefly discuss below. 

\subsection{Comparison with reconnection acceleration in turbulence}

Turbulence in magnetized fluids and reconnection are intrinsically interrelated, with reconnection being an essential part of the MHD cascade (Lazarian et al. 2020). In strong Alfvenic turbulence, the changes in the magnetic field length can be associated with the processes of reconnection and dynamo. However, in the previous sub-section, we have shown that for sub-Alfvetic turbulence, the effects of these changes are subdominant compared to the acceleration effect of the slow modes. This is especially evident if a particle ballistically samples the magnetic field of Alfvenic turbulence. Following the line of thinking in Cho \& Lazarian (2006) one can conclude that the effect is very much suppressed while introducing the scattering to enhance the particle-turbulence interactions will result in the domination of the scattering stochastic acceleration. 

The bouncing diffusion can increase the interaction of particles with magnetic turbulence. In this paper, we showed that the acceleration in incompressible sub-Alfvenic turbulence is dominated by slow modes and therefore, the effects of magnetic reconnection are subdominant.  However, we also note that for the bouncing diffusion, the acceleration in the parallel direction is bad news, as this effect increases $\mu$ and untraps the particle.

The reconnection can be an important acceleration process in superAlfvenic turbulence. The corresponding mechanism was discussed in Brunetti \& Lazarian (2011) in application to the acceleration in the environment of galaxy clusters. 

\subsection{Comparison with Transient Time Damping Acceleration}

The interactions with the magnetic mirrors might resemble the process of Transient Time Damping (TTD) acceleration (see Schlickeiser 2006, Xu \& Lazarian 2019). However, the introduced mechanism relies on different physics. For instance, the TTD acceleration require the phase speed of the compressible perturbations to be equal to the particle speed, i.e. ..... In other words, TTD happens as a particle serfs over moving magnetic compressions. On the contrary, the bouncing acceleration induced does not require any resonance interactions. The only requirement is that the particle has the pitch angle less than $\mu_c$ that allows the reflection from the magnetic bottle. 

Another difference is that the bouncing acceleration stochastically increases the perpendicular component of the particle momentum, while the TTD stochastically increases the parallel component of the particle momentum. In fact, TTD action can increase pitch angle $\mu$ removing the particles from the bouncing regime, thus shutting down the acceleration. At the same time, the TTD can be shut down if the momenta of accelerated particles is closely aligned with magnetic field. In that sense the presence of bouncing acceleration can improve the efficiency of the TTD acceleration. The bouncing acceleration does not shut up due to decreasing of the pitch angle $\mu$ of the bouncing particles.\footnote{The caveat is that the particles should not be concentrated within the range $\mu<\mu_{min}$, which can reduce the acceleration efficiency, but not shut it down.}

The two mechanisms are similar in the sense of producing highly anisotropic distribution of accelerated particles. The anisotropy in energetic particle distribution can induce instabilities, e.g. gyroresonance instability (Lazarian \& Beresnyak 2007) that will isotropize the distribution. 


\end{comment}



%\section{Discussion}

%\subsection{Physical picture of the new acceleration mechanism}



%The particles trapped in classical magnetic bottles cannot get accelerated in the stochastic second-order Fermi way if they do not experience scattering. This is a natural consequence of the preservation of the adiabatic invariant. Namely, if a particle gets accelerated in the compression of the magnetic field, it gets decelerated as the magnetic field gets to its original value. To experience second-order Fermi acceleration, a particle must experience pitch angle scattering. The latter process introduces randomness and transfers the momentum between $p_\bot$ and $p_\|$. As a result, the adiabatic invariant is not conserved, and the energy can transfer from the magnetic field oscillating around its mean value to particles. Scattering can be viewed as a process during which a particle "hides" some momentum gained in the process of adiabatic heating in the form of the momentum component that is not subject to subsequent adiabatic cooling. However, Cho \& Lazarian (2006) showed that in many cases, the acceleration from pitch-angle scattering dominates the second-order Fermic stochastic acceleration from large-scale magnetic field compressions. 


%On top of this, it is essential to note that in an incompressible fluid, the particles trapped in magnetic bottles cannot be accelerated by large-scale magnetic field compressions. This is the consequence of the simultaneous action of two adiabatic invariants. As a result, increasing the perpendicular to the magnetic field momentum is compensated by the corresponding decrease of the parallel to the magnetic field momentum. The scattering ensures that the increase of the total particle energy vanishes (see Eq. (\ref{Etot})). 


%The mechanism discussed in the paper is different. It appeals to the superdiffusion of particles perpendicular to the magnetic field that arises from the stochasticity of magnetic field lines in MHD turbulence (Lazarian \& Vishniac 1999, Yan \& Lazarian 2008, Eyink et al. 2011, 2013, Lazarian \& Yan 2014). As a result, the particles moving backward do not retrace the same magnetic field lines but instead sample different eddies. In other words, the analogy with an oscillating magnetic bottle is no more applicable. Suppose a particle is reflected by a magnetic mirror back. In that case, it samples a different magnetic bottle whose dynamics are uncorrelated with the dynamics of the compression/expansion state of the original magnetic bottle. This is the essence of the "bouncing diffusion" process introduced in LX21. This process acts on a subset of particles with sufficiently small pitch angles so that they are mirrored by magnetic fluctuations. In this paper, we considered the consequence of the bouncing diffusion on stochastic particle acceleration.

%We found that the bouncing diffusion induces efficient particle acceleration in the incompressible limit. This new acceleration process decreases the particle's pitch angles, increasing the degree of anisotropy. This guarantees that the acceleration is self-sustainable and the accelerated particles do not escape the domain of bouncing diffusion. The process is driven by the magnetic field fluctuations induced by slow modes, which in the incompressible limit are frequently termed pseudo-Alfven modes. The properties of these modes are well described in the literature and tested numerically (GS95, LV99, Cho \& Vishniac 2000, Lithwick \& Goldreich 2001, Cho \& Lazarian 2002, 2003, Kowal \& Lazarian 2010). 

%We considered the interaction of energetic particles with slow modes in the inertial range. We found that the acceleration happens in two regimes that are determined by the rate of parallel propagation of particles within the eddies. If the particle bouncing diffusion within the eddy is faster than the evolution of the slow mode eddy, this is the FPD regime. In the opposite case, the particles are in a slow propagation regime or in the SPD regime. For most of the pitch angles, the particles experience both FPD and SPD acceleration through interaction with eddies of different sizes. The contribution of two processes to the resulting stochastic acceleration is approximately equal and adds up. 

%The importance of the introduced process of acceleration is that it shows that efficient non-resonant acceleration of particles takes in the limit of incompressible MHD turbulence. This changes a lot our understanding of acceleration in various Astrophysical environments. For instance, in galaxy clusters, the medium is high better, and the incompressible approximation is valid. The acceleration of energetic particles there is usually considered versus resonance scattering introduced by fast modes (Brunetti \& Lazarian 2007). The new process alters the picture. The corresponding consequence will be elaborated on elsewhere.

%\subsection{Bouncing propagation and acceleration}

%The bouncing acceleration introduced in this paper relies on the bouncing diffusion introduced in LX21. For a number astrophysical settings discussed in LX21, e.g. for the galactic interstellar medium, latter diffusion is slower than the diffusion induces by pitch angle scattering. This can result in the relative depletion of the particle density in the bouncing domain of small $\mu$. The bouncing acceleration decreases $\mu$ and thus preventing particles from leaving the bouncing regime.



\section{Conclusions}

In a weakly compressible medium, 
both mirror trapping and pitch-angle scattering cannot lead to acceleration of CRs by magnetic compressions. 
In the former case with a 
particle trapped within a single magnetic bottle with 
oscillating magnetic compression and expansion and conserved adiabatic invariants, the stochastic acceleration does not happen.
In the latter case with efficient scattering, in each eddy a particle 
undergoes fast pitch-angle isotropization 
and 
samples both fluid compression and expansion 
that happen simultaneously in different directions, and thus its energy gain and loss cancel out. 
Moreover, in realistic MHD turbulence, 
mirror trapping does not happen, and the 
scattering by anisotropic Alfv\'{e}n and slow modes are inefficient. 



As a new diffusion mechanism,  
the mirror diffusion in MHD turbulence
identified by LX21
does not cause trapping of particles or pitch-angle isotropization. 
It takes place due to the perpendicular superdiffusion of turbulent magnetic fields regulated by Alfv\'{e}n modes 
and mirroring by slow modes. It can effectively confine CRs via stochastic mirror reflection in the direction parallel to the magnetic field. 
Under the consideration of mirror diffusion, 
% conditions: confined but not trapped, no pitch angle isotropization
the stochastic non-resonant interaction with slow-mode eddies with magnetic compression/expansion results in the mirror acceleration. 
The mirror acceleration takes place 
irrespective of the compressibility of gas. 


 







Among the slow-mode eddies with different parallel sizes along the turbulent energy cascade, we find that the ones with their lifetime comparable to the mirror diffusion timescale dominate the mirror acceleration. 
It follows that the resulting momentum diffusion coefficient does not depend on the mirror diffusion coefficient. 
The acceleration time only depends on the Alfv\'{e}n crossing time and relative magnetic fluctuation of slow modes at the driving scale of turbulence. 
In comparison with the inefficient scattering acceleration associated with inefficient scattering diffusion 
\citep{CL06}, 
the mirror acceleration serves as an efficient acceleration mechanism in a high-$\beta$ medium. 



The mirror acceleration causes stochastic increase of CR perpendicular momentum and pitch angle, and thus the acceleration can be self-sustained with the condition for mirroring satisfied. 
The mirror acceleration in a weakly compressible high-$\beta$ medium
can be applied to studying CR re-acceleration in the high-$\beta$ intracluster medium. 
This application will be investigated in our future study. 




%In the paper above, we considered particles' interaction with incompressible MHD turbulence fluctuations. We found
%\begin{enumerate}
%    \item The bouncing of particles from magnetic mirrors induces efficient stochastic acceleration by slow modes.
%    \item The new way of acceleration is identified in the process of incompressible turbulence interacting with bouncing particles.
%    \item The interaction of the particles with the turbulence cascade depends on the relation between the rate of diffusion and the fluctuation time scale of slow modes. This relation is different for different size eddies of the slow mode cascade.
 %   \item The contributions to the acceleration arising from the particles fast diffusing through the eddies and slowly diffusing through the eddies is comparable for a realistically extended cascade of slow mode turbulence.
%    \item The process of stochastic acceleration by large scale slow modes is more efficient compared to the reconnection-induced acceleration in sub-Alfvenic turbulence. It can also be more efficient than the pitch angle acceleration, especially in the presence of turbulence damping.
    
%\end{enumerate}















 







    

 

\acknowledgments
 A.L. acknowledges the support of the NSF grants AST1715754, 1816234 and NASA ATP AAH7546.   The  Flatiron  Institute  is supported by the Simons Foundation.
S.X. acknowledges the support for 
this work provided by NASA through the NASA Hubble Fellowship grant \# HST-HF2-51473.001-A awarded by the Space Telescope Science Institute, which is operated by the Association of Universities for Research in Astronomy, Incorporated, under NASA contract NAS5-26555. 
S.X. also acknowledges the support from the Institute for Advanced Study.


 
 \begin{comment}
 \appendix

 
 
 \section{Gyroresonance acceleration of slow moving particles}
 
 The gyroresonance takes place when the frequency of the incoming wave in the system of reference of a particle is equal to the gyrofrequency of the particle. The transfer of energy in the interaction is
 \begin{equation}
 
\end{equation}
As we discussed in the main text, in strong MHD turbulence the fast moving CR samples different eddies. The situation is different for a slowly moving particle. In the limiting case of very slow motion, the necessary condition for the random walk in momentum to occur is that the randomization in the pitch angle takes place. The corresponding action takes place over time determined by the diffusion of $\mu$:
\begin{equation}
D_{\mu, \mu}^{slow}\approx
\end{equation}
which provides the time for the acceleration
\begin{equation}
T_{acc}\approx
\end{equation}
Thus while the acceleration increment is the same as for the fast moving particles, the time step of the random walk can be much longer. This makes the acceleration less efficient. 

A particle moving along magnetic field lines with the speed $V_{d}$ which is not negligible with Alfven speed $V_A$, the particle in a time 
\begin{equation}
T_{tr}\approx V_A/V_d t_{l, eddy}
\end{equation}
crosses the eddy of the size $l$ in the transverse direction and gets into a new eddy. Therefore after this time, the particle starts sampling another eddy, and this restores the random walk pattern of acceleration. In other words, in a time 
\begin{equation}
T_{sl}=min[T_{acc}, T_{tr}]
\end{equation}
the particle enters the regime of random walk in terms of its angular momentum. 

These results are relevant to the acceleration of the charged dust particles via gyroresonance interactions (Yan \& Lazarian 2003, Yan et al. 2004). Our present study shows that the acceleration of dust can initially be significantly slower than it was assumed earlier. However, if the grain get velocities comparable to the Alfven velocities, the efficient resonance interaction is restored. 


\subsubsection{SuperAlfvenic turbulence}

SuperAlfvenic turbulence introduces yet another scale for the diffusion of cosmic rays. Above the scale (see Lazarian 2006)
\begin{equation}
l_A=LM_A^{-3}
\end{equation}
the MHD turbulence is hydrodynamic-like with magnetic field being easily bent by hydrodynamic motions. The magnetic field, nevertheless, is strong
enough for imposing MHD dynamics on motions with size less than $l_A$, i.e. on the motions from the range $[l_{min}, l_A]$. The description of the propagation that we had earlier is applicable to the latter range of scales. For scales less than $l_A$ one can introduce the mean magnetic field and consider the propagation parallel and perpendicular to this mean magnetic field direction. At the scales larger than $l_A$ the CRs moving ballistically experience 3D isotropic random walk with the step $l_A$ (see Brunetti \& Lazarian 2007). 

\section{Dynamics of CR on scales less than the viscous dissipation one}
\label{sec:small_r}

The dynamics of CRs changes on the scales less than the turbulence viscous dissipation scale $l_{damp}$. At scales $l<l_{damp}$ magnetic field lines preserve their identity, but are not turbulent any more. Instead, they undergo the Lyapunov separation due to the shearing motion of the critically damped eddies at scale $l_{damp}$. 

In MHD turbulence the eddies at $l_{damp}\ll L_{i}$ are strongly anisotropic and therefore 
the scale field parallel to magnetic field scale $l_{damp, \|}$ is much larger than the perpendicular to magnetic field scale $l_{damp, \bot}$, i.e. $l_{damp,\|}\ggl_{damp,\bot}$. 
If the magnetic field lines are separated by the distance $r_0$ less than $l_{damp, \bot}$, they undergo the Lyapunov exponential separation that is induced by the shear induced mostly by the eddies at the scale $l_{damp, \bot}$.  Therefore the resulting equation
\begin{equation}
\frac{d\delta}{ds}= \frac{\delta}{l_{damp, \|}}
\end{equation}
has the solution
\begin{equation}
\delta=r_0 \exp \left(\frac{s}{l_{damp, \|}}\right)
\end{equation}
For CRs to follow the divergent magnetic field lines the initial distance $r_0$ should not be less than the Larmor radius of the particles. 
 
 When $\delta$ reaches $\lambda_{damp}$, the particles laws of magnetic field line separation change. It is easy to see that to get separated by this distance the particles should travel along magnetic field lines the distance (see Lazarian 2006):
 \begin{equation}
 s_{RR}=l_{damp, \|} \ln\left(\frac{l_{damp,\bot}}{r_0}\right)
 \label{RR}
 \end{equation}
 The scale $s_{RR}$ is analogous to the separation of magnetic field lines scale introduced by Rechester and Rosenbluth (...) in the toy model of isotropic "turbulence" with motions at a single scale. In our case, instead of this single scale we employ two the scale $l_{damp, \|}$
 and $\lambda_{damp,\bot}$ that characterize the damping of anisotropic Alfvenic turbulence. The magnetic perturbations that produce magnetic field trapping are expected to be smaller than $s_{RR}$ and thus we assume that $\lambda_{trap}\gg l_{damp,\|}$.\footnote{The Alfvenic modes are responsible for shearing motions, while slow and fast modes are responsible for trapping. There are situations, e.g. in partially ionized gas, when slow modes can protrude to smaller scales compared to the Alfven modes. We do not consider such special cases in this paper.}
 
 It is evident from Eq. (\ref{RR}) that after passing a distance of several $\lambda_{damp,\|}$ the particles separate over the distance $\lambda_{damp, \bot}$ starting from which the separation of magnetic field lines prescribed by LV99 field wandering expression. 
 
 The magnetic fields in turbulent fluid are separating from each other at an accelerated rate given by Eq.(\ref{perpen}). This accelerated separation arises is due to the fact that the larger the separation the faster the magnetic field lines diverge. This effect can be viewed as the gradual increase of the dispersion of the magnetic field line directions that determine the motion of the particles as the initial separation between the particles increases. If CRs follow magnetic field lines without scattering or trapping, after propagating a few distances of $l_{damp,\|}$ the CRs get the regime of LV99 superdiffusion. 
 
 There exist, however, a possibility that CRs can be scattered on the scales much less than $l_{damp,\|}$. For instance, streaming instability can induce resonance fluctuations that make the CRs mean free path $\lambda$ less than $l_{damp, \|}$. In this case, 
 the retracing the particles is possible over limited scales and the step in the perpendicular direction is $l_{damp, \perp}/l_{damp,\|}$. To get the $l_{damp,\bot}$ a CR has to pass a distance give by $s_{RR}$. 
 

 
  
 
 

\section{Boltzmann equations for CR propagation}

As we discussed in the previous section, there are two separate populations of cosmic rays that participate in the process of the parallel CR propagation, one consists of the trapped particles and has the distribution function 
$f_{t}$, while the other is of the freely scattered particles with the distribution function $f_u$. The two populations exchange particles, the transfer from $f_t$ to $f_u$ has the rate $\nu_{tu}$ and the transfer back has the rate
$\nu_{ut}$. To describe the process one can write two coupled Boltzmann equations in an analogy with the accepted way of treating the particle transport (Earl 1974):
\begin{equation}
 \frac{\partial f_u}{\partial t} +\mu V_u \frac{\partial f_u}{\partial z}=\frac{1}{2} \frac{\partial}{\partial \mu} \phi \frac{\partial f_u}{\partial \mu} - \nu_{ut}^{-1} f_u+ \nu_{tu} f_{t}
 \end{equation}
 \begin{equation}
 \frac{\partial f_t}{\partial t} +\mu V_t \frac{\partial f_t}{\partial z}=\frac{1}{2} \frac{\partial}{\partial \mu} \phi \frac{\partial f_t}{\partial \mu} - \nu_{tu}^{-1} f_t+ \nu_{ut} f_{u}
 \end{equation}
 where $\phi$ is determined by the properties of realistic anisotropic MHD turbulence. The exchange time-scale from the propagating state to the trapped state $\nu_{ut}^{-1}$ was calculated in Xu \& Lazarian (2020):
 \begin{equation}
 \nu_{ut}^{-1}\approx \frac{28}{5\pi} \frac{B_0^2}{\delta B_f^2} \left(\frac{v}{L\Omega}\right)^{-1/2} \frac{v^2}{\Omega}\left[4-\sqrt{\mu_{cf, f}} (5-\mu^2_{cf,f})\right]
 \end{equation}
 the return time scale is
 \begin{equation}
 \nu_{tu}^{-1}\approx 
 \end{equation}
 
 
 



\section{Trapping and Second Order Fermi acceleration}
\label{B}

The trapping induces a new type of acceleration that we terms Trapping Acceleration (TA). The physics of TA is very simple. 
In turbulent environments the trapping is due to magnetic perturbations that are moving along magnetic field lines. Consider, for instance, the trapping induced by slow modes. Such perturbations move along the magnetic field with the velocity of $V_{sl}$. For considering purely compressible perturbations it is natural to assume that we are dealing with the magnetically dominated low $\beta$ medium in which slow modes propagate with the sound velocity, i.e. 
$V_{sl}=c_s$. For high $\beta$ media slow modes are compressions of magnetic field in essentially incompressible media and their effects on CR acceleration is different (see \S \ref{A}). 

On interacting with the corresponding magnetic bottle the particle momentum changes by 
\begin{equation}
\frac{\Delta p}{p}\approx \frac{V_{sl}}{v_{par} \mu},
\end{equation}
where $v_{CR, \|}$ is the parallel to magnetic field CR velocity. This velocity changes for the trapped particles, so one can use the mean parallel velocity for $v_{part}$ for our estimates of second order Fermi acceleration. 
As a result the momentum diffusion coefficient for a given $\mu<\mu_{crit}$ is
\begin{equation}
D_{p}({\mu})=\frac{(\Delta  p)^2}{\Delta t}\approx p^2\left(\frac{V_{sl}}{v_{CR}\mu}\right)^2 v_{CR} \mu k_{\|},
\label{D_pmu}
\end{equation}
where it is taken into account that the time between collisions with moving magnetic mirrors is $1/(v_{par}\mu k_{\|}$, where $k_{\|}^{-1}$ is the trapping scale corresponding to the particle with a given $\mu$ (see Eqs. (\ref{bkslow}) and (\ref{kparfast})). To find the total momentum diffusion coefficient Eq. (\ref{D_pmu}) should be integrated over $\mu$ for the range of the trapping angles:
\begin{equation}
D_p=\int_0^{\mu(crit)} D_p(\mu)\frac{d\mu}{\mu}
\end{equation}

The treatment above is applicable for compressible media, e.g. to the slow modes in low $\beta$ plasmas.
For the slow modes one gets with the help of Eq. (\ref{bkslow})
\begin{equation}
D_{p,slow}=\frac{1}{5\mu^5_{cr, s}} \frac{v_{CR}}{L} p^2 \left(\frac{V_{sl}}{v_{CR}}\right)^2 \left(\frac{\delta B_{sl}}{B_0}\right)^2,
\end{equation}
where $V_sl$ is the phase velocity of slow modes, which is the   the sound velocity $c_s$ in the low $\beta$ plasmas. For most cases $v_{CR}$ is the speed of light $c$. 

For the fast modes, similarly, using Eq. (\ref{kparfast}) it is easy to get
\begin{equation}
D_{p,fast}=\frac{1}{9\pi \mu^9} \frac{v_{CR}}{L} p^2 \left(\frac{V_{f}}{v_{CR}}\right)^2 \left(\frac{\delta B_{sl}}{B_0}\right)^4,
\end{equation}
where $V_{f}$ is the velocity of the fast wave, which and $c_s$ for $\beta \gg 1$ and the integration over the isotropic distribution of fast wave is performed.  


Note that the random walk arguments to the situation when particles bounce between magnetic mirrors as every time the particles interact with a new mirror. Indeed, the LV99 dispersion provides the perpendicular displacement of $l_{\bot}$ while the particle moves along the magnetic field lines the distance $l_{\|}$, where $l_{\bot}$ and $l_{\|}$ the perpendicular and parallel scales of the magnetic eddy. In other words, a particle does not bounce back twice within the same eddy. As the eddies are not correlated, the CR gets its increments of velocity in a random walk fashion. 

The process of acceleration produced by this mechanism provides an anisotropic acceleration that favors the increase of the parallel component of particle momentum. The increase of the parallel momentum of the particle induces the "un-trapping" of particles. For a particle with the pitch angle $\nu$ to get into the escape cone the change of the angular momentum 
\begin{equation}
(\delta p)^2\approx p^2 (\nu_{crit}-\nu)^2
\end{equation}
This requires time
\begin{equation}
t_{ac}\approx (\delta p)^2/D_p\approx \frac{(\nu_{crit}-\nu)^2}{\left(\frac{V_{sl}}{v_{CR}}\right)^2 \frac{v_{CR}}{\Delta_{trap}}}
\label{tac}
\end{equation}
where the expression for the momentum diffusion coefficient given by Eq. (\ref{D_pmu}) was used. In other words, the time scale $t_{ac}$ presents another time scale for the particles to leave the trap. Compared to the 
removal of trapping by scattering, this mechanism of particle escape relies only on the compressible fluctuations that initially trap particles (see Xu et al. 2017). In some settings, e.g. in turbulence in partially ionized gas, for a range of scales the slow modes may be the dominant scattering agent. 

The removal of trapping due to acceleration given by Eq. (\ref{tac}) should be compared to the time of $\mu$ randomization due to cosmic ray scattering,
$1/\Gamma_s$ 
given by Eq. (\ref{}). 

One can compare the efficiency of TA and the TTD acceleration by slow modes. The corresponding expression for the TTD is given in Xu \& Lazarian (2019). ...

We note, that that if the scattering is not sufficiently efficient, the anisotropic distribution of CRs induces gyroresonance instability that isotropises CRs.

In the presence of perpendicular diffusion mediated by trapping, acceleration is also possible in terms of the perpendicular momentum of CRs. In Appendix \ref{A}, we discuss the betatron acceleration of the CRs perpendicular momentum that is possible even in the case of an incompressible medium. The earlier study of this process in Cho \& Lazarian (2006) implicitly assumed that the only way for the CR to be returned back is through scattering. In incompressible fluid, this results in suppression of the acceleration, however.






\subsection{Physical picture of acceleration}
 
 Naturally, the choice of the sign of $\frac{v_{l,\|}}{l_{\|}}$ is varying and the time derivative of the CR kinetic energy is changing its sign. Evidently
 if 
 \begin{equation}
 E^{kin}_\bot=2E^{kin}_\| 
 \label{equal}
\end{equation} 
 the change of kinetic energy is zero. However, this equality is not self-sustained, unless there is a strong isotropization
 process at work. Indeed, if the condition given by Eq. (\ref{equal}) is fulfilled at one moment, the kinetic energies 
 corresponding to the parallel and perpendicular motions are changing according to Eq. (\ref{E1}) and (\ref{E2}) and this 
 violates the equality in Eq. (\ref{equal}). As a result, depending on the signs of 
 \begin{equation}
 I_1=\frac{v_{l,\|}}{l_{\|}}
 \label{I1}
 \end{equation}
 and 
 \begin{equation}
 I_2=E^{kin}_\bot-2E^{kin}_\|
 \label{I2}
 \end{equation}
 the energy in the CR is going to increase or decrease. 
 
 To understand the dynamics of the acceleration/deceleration we can consider two eddies, N1 having $I_1>0$, and eddy N2 having $I_1<0$. Assume that initially the equality given by Eq. (\ref{equal}) is fulfilled. 
 For N1 eddy, according to Eq. (\ref{E1}) and (\ref{E2}) $\dot{E}^{kin}_\bot>0$ and $\dot{E}^{kin}_\|<0$. Therefore the kinetic energy of perpendicular motion is increasing, while the one corresponding to the parallel motion is decreasing. As a result $I_2$ given by Eq. (\ref{I2}) gets positive. For both $I_1$ and $I_2$ that are $>0$, Eq. (\ref{Etot}) testifies that the total energy $E^{kin}$ is increasing and the increase of the energy gain is increasing in time as $I_2$ increasing. Taking into account that the formal solution of Eq.(\ref{E1}) and (\ref{E2}) are, respectively,
 \begin{equation}
E^{kin}_\bot=E^{kin}_{\bot, 0} \exp{\left[\frac{t}{l_\|/v_\|}\right]}
\end{equation}
and  
\begin{equation}
E^{kin}_\|=E^{kin}_{\|, 0} \exp{\left[\frac{-2 t}{l_\|/v_{\|}}\right]}
 \end{equation}
 where $E^{kin}_{\bot, 0}$ and $E^{kin}_{\|, 0}$ correspond to the initial condition of $I_2=0$, we see the reason why the total energy is increasing. For 
 $t\ll l_\|/v_\|$ the increase is proportional to $(E^{kin}_{\bot, 0}+ 2E^{kin}_{\|, 0}) t/(l_\|/v_{\|})$. 
 
 For the N2 eddy, $I_1<0$ and the kinetic energy associated with the parallel motion is increasing according Eq. (\ref{E2}), while the perpendicular component is decreasing as it is seen from Eq. (\ref{E1}). As a result, for $I_2=0$ initially, at later times $I_2$ gets negative. For $t\ll |l_\|/v_\||$ the decrease is proportional to $-(E^{kin}_{\bot, 0}+ 2E^{kin}_{\|, 0}) t/|l_\|/v_{\|}|$,  where we take the modulus of $l_\|/v_\|$ accounting for the fact that $v_\|<0$. Thus over small time spans the process of the increase of the energy in the eddies is a random walk.

 





\subsection{Dynamics of pitch angle}

To find how the acceleration changes $\mu$ it is useful to consider the change of square of the tangent of the pitch angle:
\begin{equation}
\tan^2\gamma (t)=\frac{P_{\bot, t}^2 }{P_{\|, t}^2}=\frac{(P_\bot +\delta p_\bot)^2}{(P_\| +\delta p_\|)^2}
\label{tan}
\end{equation}
where $P_\bot$ and $P_\|$ are initial values of the CR perpendicular and parallel . 
For $\delta p \ll P$, one can expand the expression in Eq. (\ref{tan}) keeping only the quadratic $\delta p/P$ terms. Then, providing averaging one can after simple algebra:
\begin{equation}
\langle \tan^2 \gamma (t) \rangle \approx \tan^2\gamma_0 \left( 1+\frac{\langle (\delta p_\bot)^2 \rangle}{P_\bot^2}-\frac{\langle (\delta p_\|)^2\rangle}{P_\|^2} -2\frac{\langle \delta p_\| p_\bot \rangle}{P_\| P_\bot} \right),
\label{tant}
\end{equation}
where $\langle... \rangle$ denote averaging took into account that $\langle p_\bot \rangle=0$ and $\langle p_\| \rangle=0$.  
 The increments of parallel and perpendicular momenta given by Eq. (\ref{E1}) and (\ref{E2}) and those can be written in the form:
\begin{equation}
\delta p_\bot=\frac{1}{2} p_{\bot} \xi,
\end{equation}
and
\begin{equation}
\delta p_\|=- p_\| \xi,
\end{equation}
where $\xi$ is is proportional to the eddy velocity $v_\|$, which for the problem at hand is a random variable. Substituting those into Eq. (\ref{tant}) one gets
\begin{equation}
\langle \tan^2 \gamma\rangle (t)\approx \tan^2\gamma_0 (1+\fract{1}{2}\xi^2),
\label{tanf}
\end{equation}
where 
\begin{equation}
\xi^2=v_\| m P_\| \frac{v_\|}{l_\|} t
\end{equation}


Eq. (\ref{tanf}) reflects the dynamics of the pitch angle for a given $P_\bot$ 
and $P_\|$. Using the trigonometric identity $\mu^2=1/(\tan^2\gamma +1)$ it is easy to see that
\begin{equation}
\mu^2 \approx \mu_0^2 \left(1-\frac{1}{2}\xi^2(1-\mu_0^2)\right),
\end{equation}
which describes the dynamics of the pitch angle resulting form the Trapping Acceleration (TA). 

The change of $\mu$ due to the TA can be estimated as:
\begin{equation}
G_{TA}=\frac{1}{\mu^2}\frac{d\mu^2}{t} \approx  -\frac{1}{2} v_l m P_\| \frac{v_\|}{l_\|} (1-\mu^2).
\end{equation}
The decrease of $\mu$ induces stronger trapping. In these circumstances for those $\mu$ that make absolute value of $G_{TA}$ larger than the 
the rage of $\mu$ randomization via resonance scattering (see Eq. (\ref{}) the CRs stay trapped. By equating these two rates of $\mu$ change one
can define the critical cosine of pitch angle:
\begin{equation}
\mu_{cr, TA}\approx
\end{equation}
For $\mu<\mu_{cr,TA}$ the CRs stay trapped increasing their perpendicular momentum. This can result in gyroresonance instability, eventially (see Kulsrud ..., Lazarian \& Beresnyak 2006). 










 \subsection{Resonance scattering}
 
Resonance scattering is a process known to randomize CRs angular momenta. If the scattering time $\lambda_{mfp}/v_{pt}$ is much less than $|l_\|/v_\||$ then the increase of the energy is going to be a random walk with the diffusion coefficient:
\begin{equation}
D_{E, add}\approx  (E^{kin}_{\bot}+ 2E^{kin}_{\|})^2 \frac{\lambda_{mfp}}{v_{pt}} \left(\frac{l_\|}{v_\|}\right)^2,
\end{equation}
 provides the increase of energy which is subdominant to the increase of energy that arises from the resonance scattering, i.e. 
\begin{equation}
D_{E, res}\approx E^{kin}^2 \left(\frac{V}{v_{pt}}\right)^2 \frac{v_{pt}}{\lamda_{mfp}},
\end{equation}
where the $V$ is the velocity of the perturbations that induce scattering. Indeed, it is easy to see that for $D_{E, add}\approx (\lambda_{mfp}/l_\|)^2 (v_\|/V)^2D_{E, res}$, provided that $\lambda_{mfp}\ll l_{\|}$. 

For typical incompressible MHD turbulence driven at large scales we can expect the $\lambda_{mfp}\gg l_\|$. In this situation the momentum distribution of CRs can be very anisotropic.  Consider eddy N1 with $I_1>0$ and $I_2>0$. Then the total energy is increasing and according to the Eq. (\ref{Etot}) the energy is increasing thus increasing $I_2$.  For the same eddy, if $I_2<0$, the total energy is decreasing, but the absolute value of $I_2$ is decreasing, which slows down the decrease of the total energy. 

If one considers eddy N2 with $I_1<0$, then for $I_2>0$ the total kinetic energy is decreasing according to Eq. (\ref{Etot}), the value of $I_2$ is also decreasing, which slows down the rate of energy decrease. For $I_2<0$, however, the total kinetic energy is increasing together with the absolute value of $I_2$. This result in the accelerated increase of the energy. 

The two cases above suggest that for some initial partitions between the kinetic energy of CRs parallel and perpendicular motions we may have the increase of energy which is larger than its decrease. This would result in the anisotropic distributions for which the further changes of energy would be different. For instance, if as the result of the action of eddy N1 all CRs get into $I_2>0$, the action of the eddy N2 will be to decrease the total energy.

In the absence of trapping, the acceleration can happen only on the time of ballistic crossing of the eddy. 
This case that was considered in Cho \& Lazarian (2006)  corresponds to acceleration of the CR as it propagates over the distance of the order of the mean free path $\lambda_{\|}$. For this propagation Cho \& Lazarian (2006) considered the interaction of the CR with the magnetic field of the eddy with $l_{\|}=\lambda_{\|}$. The resulting acceleration rate was found to be relatively small. 


The situation changes if the trapping is taken into account. As the gradients of the parallel velocity amplitude, i.e. 
$|v_\|/l_\||$ in the turbulent fluid increases with the decrease of the scale, it is natural to assume that for a CR with a cosine of pitch angle $\mu$ 
the most important for the acceleration will be the fluctuations with the maximal parallel wave number corresponding to the trapping. For the slow modes the corresponding wavenumber is given by Eq. (\ref{kparslow}). 





\section{Trapping and First Order Fermi acceleration}


\subsection{Trapping and acceleration in shocks}

The acceleration in shocks is the first example of the First Order Fermi acceleration. The process in its simplest form is due to particles bouncing back and forth from the upstream and downstream regions of the shock (Longhair 2011). The shock itself creates in this picture a zone of contraction with the postshock gas moving with 3/4 of the shock velocity towards the stationary medium in front of the shock. As CRs do not feel the shock itself, the acceleration is easy to consider in the frame of reference moving at the 3/8 of the velocity of shock. In this frame of the observer sees the postshock gas moving towards him with 3/8 of the velocity from behind and the preshock gas moving with 3/8 of the velocity towards him from the front. Naturally, in this reference frame the acceleration of CR happening in every "collision" with postshock and preshock matter. 

In the traditional Diffuse Shock Acceleration (Krymsky..., Bell 1978, Blandford \& Ostriker 1978) the "collisions" happen due to resonance scattering. The reversing the direction of the particle is a relatively show process as pointed out in 
Jokipii (....). There it is shown that some observational data, e.g. ....,  is impossible to explain using this process. The magnetic mirrors can reflect CRs faster and this increases the rates of CR acceleration in shocks. The increase of magnetic field intensity can arise from magnetic field generation in the pre-shock and post-shock areas due to the dynamo action. For instance, Balsara ... considered the dynamo in the post shock gas while Beresnyak et al. (...) and Xu \& Lazarian (...) considered dynamo in the precursor of the shock. 

For shocks the differences in acceleration in the parallel and perpendicular shocks have been considered. For instance, Jokipii .... championed the high efficiency of the perpendicular shocks for the acceleration. He pointed out that if the CR is moving ballistically along a magnetic field line while the turbulent line crosses the shock many times, the acceleration will not depend on the slow processes of diffusion that in the diffusive shock acceleration (DSA) model return the CR to the region of acceleration. It was shown in Lazarian \& Yan (2014), however, that this scenario contradicts to what we know about the properties of Alfvenic turbulence. Indeed, the magnetic field wondering over the inertial range of turbulence induces the dispersion of magnetic field lines that increases as $s^{3/2}$, where $s$ is the distance tracked along the magnetic field line. As a result, the probability of the particle returning back to the shock gets low. At the same time LY14 discussed a possibility that the dynamo in the precursor of the shock can produce a ribbon of high intensity irregular magnetic field with the scale of the magnetic eddies less than the thickness of the shock precursor. The process of precursor dynamo was discussed in Beresnyak et al... with some properties of it tested in ....Santos-Lima \& Lazarian .... Provided that the CRs stay appreciable time in the ribbon, the process analogous to that suggested by Jokipii ... is feasible, as on the scales larger than $l_A$ the magnetic trajectories are diffusive and, as a result, multiple crossing of the shock by CRs moving ballistically along magnetic field lines is possible. CR trapping on the scale $l_A$ is expected to contribute to extending the time for the CR to stay within the ribbon, thus increasing the efficiency of this ballistic shock acceleration (BSA) process. 

An additional boost of the acceleration efficiency for the BSA can be obtained for high energy CRs, that have the Larmor radii larger than the shock front determined by the media mean free path. If such particles propagate with small $\mu$ along the magnetic field line that crosses the shock at a small angle,
the CR can get the boost of their momentum. Consider a magnetic field of strength $B$  at an angle $\theta$ to a shock. As a CR with the momentum $p$ and the Larmour radius $r_L$  gyrates in this magnetic field it passes the distance $l_{el} \approx r_L \cos\theta/\sin\theta$ while crossing the shock and interacting with both the preshock and postshock media in the process of the gyration. The number of gyrations can be estimated as
\begin{equation}
N\approx \frac{\cos\theta (1-\mu^2)^{1/2}}{\mu \sin\theta}.
\end{equation}
 while crossing the shock and interacting with converging magnetic field the CR gets an increase of the momentum $V_{sh}/p(1-\mu^2)^{1/2}$, then the increase of the angular momentum per the ballistic crossing of the shock is
 \begin{equation}
 \delta P_{sh} \approx \frac{V_{sh}}{p} \frac{\cos\theta}{\mu \sin\theta},
 \end{equation}
 which for small $\theta$ and $\mu$ can significantly enhance the efficiency of shock acceleration. If CR experience trapping on the distance $l_{el}$, this is expected to increase 
 $\delta P_{sh}$ further. It is important to note that the process that we discuss increases the perpendicular component of the angular momentum thus
 decreasing $\mu$. The latter further enhances trapping according to Eq. (\ref{mubk}). 
 
 Bouncing back and forth between magnetic mirrors can also make the acceleration in terms of the parallel momentum in a shock more efficient by transferring it to the regime of the BSA. This scenario can be realized in the case of trapping of CR in the subAlfvenic turbulence by the large scale trap at
 $L_i$ scale. As we discussed earlier, the perpendicular propagation over the scales larger than $l_A$ is a diffusion with the time step $\approx L_i/v_{pt}\mu$ and the perpendicular transposition during this step $\sim l_A$. Assuming that the energy injection of compressible motions is isotropic we find that it takes
 $(L_i/l_A)^3=M_A^4$ steps to leave trap at the scale $L_i$. During this time the increase of the CR momentum is going as
 \begin{equation}
 \frac{dp_\|}{dt}\approx p \frac{V_{sh}}{v_{pt}\mu},
 \end{equation}
 i.e. the final momentum of the CR can be 
 (not understood commented)
 %$MA_A^4 p (V_{sh}/(v_{pt\mu)$. 
 
 
 Another case for the BSA is presented for the perpendicular shock propagating in inhomogeneous media if the shock encounters a turbulent magnetized 
 region, e.g. turbulent flux tube parallel to the shock front. The rate of the diffusion of the CRs out of the flux tube is limited by the Bohm diffusion, while perpendicular to the tube direction CR can move superdiffisively due to the magnetic field meandering (see \S...). For an idealized situation of a magnetic flux tube with $M_A<1$ turbulence. As we discussed earlier, the mean speed of the particle perpendicular to the magnetic field is $v_{pt}(1-\mu^2)^{1/2} M_A$. If the shock speed is $V_{sh}$ and the tube thickness is $l_{tube}$, then the acceleration time is $l_{tube}/V_{sh}$ and the particle can cross the shock. After simple algebra for sufficiently small $M_A$ one can get that the increase of the 
 particle momentum after crossing such a tube does not depend on the thickness of the tube and can be $\delta p\sim p M_A$. 
 
 
 \subsection{Trapping and acceleration in reconnection regions}

The First order Fermi acceleration was proposed in DeGouveia dal Pino \& Lazarian 2005 and Lazarian (2005) as a part of the model of LV99 turbulent reconnection. Figure 1 reproduced from Lazarian (2005) illustrates the process. 
There the particles bounce back and force from the magnetic bottles created by the base of the reconnection regions. As a result of this bouncing the CR is constantly gaining increments of the momentum parallel to magnetic field component of momentum with every bouncing. This is a particular way of acceleration with the distribution function of CRs being anisotropic and increasing its anisotropy as a result of the acceleration. This gradually increases $\mu$ removing the trapping of CRs.

CRs with sufficiently large gyroradius $r_L$ will interact with converging magnetic fluxes increasing the perpendicular momentum and thus decreasing $\mu$ (see Lazarian ....). This can increase trapping. The difference of the effect from that we considered for the increase of CR perpendicular momentum of the shocks is that the CR interacting with two fluxes of the opposite direction does not gyrate, but is involved in the fast motion perpendicular to the direction of both fluxes as it gyrates. 

In the presence of resonance scattering the given type of acceleration is suppressed. However, the reflections induced by trapping are different, namely, they do not randomize the pitch angle. As a result, the trapped particles can spend more time experiencing the process of magnetic field shrinking and gaining energy as a result of this. 








\section{Damping of Alfven waves in MHD turbulence mediated by density fluctuations}

In Farmer \& Goldreich (2004) and Lazarian (2016) the non-linear damping of Alfven modes emitted in the process of streaming instability by the MHD cascade was considered. The idea of this damping is very simple. The emitted Alfven waves propagating along local magnetic field get the non-zero $k_\bot$ wavenumber and therefore can efficiently cascaded by the eddies with the same perpendicular wavenumber. However, the Alfven waves propagating through a media of the changing density are also gaining $k_\bot$ as the local Alfven velocity changes. While passing through a density fluctuation of the size $l_{den}$ with the density $\delta \rho_l$, the Alfven mode gets deflected by an angle $(v_1-v_2) T/\lambda\sim \delta \rho_l/(2\rho)$. T the fluctuations of density should be taken at the scale $l=\lambda$ and the effect depends on the spectrum of density in the ambient medium. Unlike velocity fluctuations, the spectrum of turbulent density is not universal and it depends both on the medium, i.e. one phase or multi-phase, and on the sonic Mach number $M_s$. For instance, the density fluctuations in high $M_s$ media have shallow density spectrum with the density fluctuations increasing with the decrease of the scale (see Beresnyak et al. 2005, Kowal et al. 2007). Such density fluctuations increase significantly the turbulent damping of streaming instability at low CR energies. 

It is shown in Lazarian (2016) that the waves are being cascaded by the motions of the eddies with perpendicular scale $x=\lambda/\sin\theta$, where $\theta$ is the angle determined by the variations of the direction of the wavefront of the propagating wave. In the case of Alfven waves being damped by the Alfvenic turbulence this angle is determined by the variations of the magnetic field direction at the scale $x$, namely $\theta \sim (\delta B_x)/B$. In the presence of density fluctuations an additional source of the propagation of Alfven waves is related to the refraction of Alfven waves. If the spectrum of the density is shallow, the maximal refraction is expected to be at the scale $x$, i.e. one should use $\theta\sim 1/2\rho_x/\rho_0$, where for the fluctuations that follow the Kolmogorov scaling:
\begin{equation}
\rho_x\approx \rho_L \left(\frac{x}{L}\right)^{1/3}
\end{equation}
The relative efficiency of wave non-linear damping depends then on whether the ratio of $(\rho_L/\rho_0)$ is larger or smaller than $M_A$. The former ratio increases with the sonic Mach number and therefore for low $\beta$ medium the damping by the new process can more important even if $M_s<1$. 

According to Cho \& Lazarian (2003) slow modes create fluctuations of densities $\delta \rho_L/\rho_0\sim M_s$ (see \S ...). The corresponding scaling for $\delta B_L/ B\sim M_A$. Therefore for $M_s/M_A\sim V_A/c_s>1$, i.e. for magnetically dominated media the damping of MHD waves by turbulence mediated by slow modes is dominant, provided that the slow modes are not suppressed through collisionless dampling. The corresponding expression for the damping is ...
\begin{equation}
\end{equation}
which provides the damping scale for slow modes cascaded by Alfvenic turbulence 
\begin{equation}
\end{equation}

Table ... summarizes the corresponding damping scales for slow modes in different media. It is clear, however, that for a range of CR energies the suppression of streaming instability mediated by slow modes can be dominant. 

For $M_s>1$ or with the multiphase fluid, the spectrum of density may be shallow (see Beresnyak et al. 2005) with the highest density perturbation being at the smallest scales. In this situation, the maximal fluctuations of density are expected at the scale $\lambda$ and therefore $\theta\approx 1/2 \delta \rho_\lambda/\rho_0$. The corresponding damping rate of the Alfven modes is 
\begin{equation}
\end{equation}
\end{comment}

\bibliographystyle{aasjournal}
\bibliography{xu}


\end{document}






\section{"strong trapping"}
\label{subd}



\subsection{Subdiffisuion}
Sufdiffusion is the process explored in many papers. The idea of the process that a CR retraces its trajectory while being scattered backwards. In the presence of the magnetic field wandering this re-tracking is highly improbable if CRs are moving along the magnetic field lines which turbulence corresponds to the regime of strong Alfvenic turbulence (Yan \& Lazarian 2008). In Lazarian \& Yan (2014) another situation corresponding to the re-emergence of subdiffusion was mentioned. We elaborate this case of sub-Alfvenic turbulence below. 

This effect is not present if particles sample Alfvenic turbulence at the scales less than $l_A$. This can happen in two cases. One is that the damping scale of turbulence $l_{damp}$ is larger than $l_A$, which can be realized, e.g. in partially ionized gas, where the ion-neutral collisions increase the scale $l_{damp}$ (Xu et al. 2016, see also Xu \& Lazarian 2017). Another possibility for the subdiffusion is that the Larmor radius of the CR $r_L>l_A$. 

Another and more artificial case of subdiffusion can potentially happen if the turbulence is 2D. Such situations can take place in disk systems for turbulence scales larger than the disk scale height $H_{disk}$. The physics of 2D turbulence is not fully understood (see discussion in Eyink et al. 2011), but we believe that the superdiffusive (Richardson-type) behavior is not applicable to this regime. As a result, for CR with $r_L>H_{disk}$ one may observe the subdiffision perpendicular to the magnetic field in the disk. 

\subsection{Strong trapping}

In the case of subdiffusive propagation the effects of trapping can have different consequences for the CR propagation. 

The discussion in Xu \& Lazarian (2020) allowed us to identify for a particle with the pitch angle $\mu$ the corresponding scale of trapping $\Delta_{trap, \mu}\sim 1/k_{\|, \mu}$. The particle can be trapped for time 
(not understood, commented)
%$t_{tr}=min{t_{eddy, k}, t_{scat, \mu}$,
where $t_{eddy}$ is the lifetime of an eddy, while $t_{scat}$ is the time of "rescuing" of the CR from the trap:
\begin{equation}
t_{trap, \mu} \approx \left[\frac{1}{\mu_{cr}-\mu} \int_{\mu}^{\mu_{cr}} \Gamma_{scat} d\mu\right]^{-1}
\end{equation}
where $\Gamma_{scat}$ is the scattering rate due to the CR interaction with turbulence. The time scale $t_{eddy, k}$ is the lifetime of the eddy with the scale $k_{\|}$ corresponding to a given $\mu$. To find the mean trapping time
and mean scale of trapping both $\Delta_{trap, \mu}$ and $t_{trap,\mu}$ should be averaged over $\mu$. 

If the mean free path of untrapped CRs is $\Delta_{scat}$ is the scattering time, i.e. the time to turn the CR momentum direction $t_{scat}$ the diffusion coefficient for the parallel diffusion of CRs is
\begin{equation}
D_{\|}=\frac{(Delta_{trap} +\Delta_{scat})^2}{t_{trap}+t_{scat}}
\label{Dl}
\end{equation}
In many cases $t_{trap}$ is expected to be much larger than $t_{scat}$. Similarly, in many cases $\Delta_{scat}\gg \Delta_{trap}$. As a results Eq. (\ref{Dl}) reduces to
\begin{equation}
D_{\|}\approx \frac{\Delta_{scat}^2}{t_{trap}}
\end{equation}

\subsection{Comparison with earlier studies}

The theory of CR propagation initially was developing on the basis of the incorrect theories of MHD turbulence. A naive assumption that CR dynamics can be treated assuming that the astrophysical magnetic field 
can be represented as a sum of strong regular magnetic field $B$ and the turbulent magnetic field $\delta b$, where $\delta b\ll B$ and the the fluctuations are isotropically distributed was initially adopted. The 
assumption $\delta b\ll B$ is at the core of the quasi-linear theory of CR propagation (see Schlikeiser 2006). To reconcile the theory with the in situ measurements of energetic particle propagation, later the model 

\section{Example: escape from eddies and CR acceleration}

{\bf This looks to be wrong. The rates of perpendicular diffusion are not so large.}
 
The finding that the perpendicular propagation depends only on the rate of magnetic field wandering and is independent on the scattering rate. This changes a lot in earlier studies of CR physics. For instance, in Chandran (2003) the acceleration of CR by large scale eddies related to slow modes was considered and two cases were identified: slow diffusion out of the eddy and a fast diffusion out of the eddy. The diffusion was assumed to be parallel to magnetic field lines. 


The resulting energy diffusion coefficient arising from the compression on the parallel scale $l_\|$ (see Appendix) is independent of the parallel diffusion coefficient (contrary to Chandran 2003) and is given by an expression:
 \begin{equation}
 D_{E,scat, l}\approx p^2 \left(\frac{v_{l}}{l_\|}\right)^2 \left(\frac{l_\|}{v_{pt}\mu}\right)=p^2 \left(\frac{v_l}{v_{pt}\mu}\right)^2 \frac{v_{pt} \mu}{l}
\end{equation}
 
 As the acceleration is taking place over the entire range of eddies from the viscous/resistive scale $l_{min}$ to the injection scale $L_i$. Using the scaling for slow modes is given by Eq. (\ref{vslow}) and using Eq. (\ref{lpar}) one gets:
 \begin{equation}
 D_{E, scat}\approx p^2 \left(\frac{v_L}{v_{pt}\mu}\right)^2 \frac{v_{pt} \mu}{L_i}M_A^{5/3} \ln\left(\frac{L_i}{l_{min}}\right),
 \end{equation}
 
 which differs significantly from the corresponding expression in Chandran (2003). 
 
 Our discussion above shows that in the MHD turbulence the acceleration takes place in the regime of the CRs fast passing through the turbulent eddies compared to the time for the evolution of these eddies. The exception is only the degenerate case of $\mu\rightarrow 0$.\footnote{The discussion of how to deal with the 90 degree scattering can be found in e.g. Yan \& Lazarian (2008) and Xu \& Lazarian (2019).} This significantly changes and simplifies the picture of turbulent acceleration (compare to Table 1 in Cho \& Lazarian 2006). 
 We note that the slow diffusion limit discussed by Ptuskin (1988) can take place for the compressions in superAlfvenic turbulence at the scale larger than $l_A=L_i M_A^{-3}$.  
 
%Ptuskin, V. 1988, Soviet Astron. Lett., 14, 255  


 
\bibliographystyle{apj.bst}
\bibliography{xu}


\end{document}